% Contribution by /root/integral_fock, GPT-6 Astra Pro; effort not exposed.
% Findings IF-01--IF-11. This file is intended for inclusion after amsmath,
% amssymb and amsthm are loaded. No external macros other than standard LaTeX.
\section{The integral Fock comparison on the actual exterior lattice}
\label{sec:integral-fock}

Put $M_n=\bigoplus_{a=1}^n(\mathbb Z e_a\oplus\mathbb Z f_a)$,
$V_n=\bigwedge^*M_n$, $p_a=e_a\wedge f_a$, and
\[
 L=\Bigl(\sum_a p_a\Bigr)\wedge-,\qquad
 D=\sum_a\iota_{f_a^\vee}\iota_{e_a^\vee}.
\]
The lowering here is ordinary contraction, not the inverse-string lowering
of HODGE\_SIGN SL-08. All lattices below are specified over $\mathbb Z$.

\subsection{IF-01: an explicit integral Fock intertwiner}
For a subset $S\subseteq\{1,\ldots,n\}$ and choices
$\epsilon_a\in\{e_a,f_a\}$ for $a\in S$, let
$\mathcal P_{S,\epsilon}$ be the submodule of
$\mathbb Z[x_a,y_a:a\notin S]$ consisting of polynomials homogeneous of
column degree one in each pair $(x_a,y_a)$. Give its monomial basis the
factorial Fock pairing. This basis is orthonormal, since every exponent is
zero or one. Set
\[
 \mathcal F_n=\bigoplus_{S,\epsilon}\mathcal P_{S,\epsilon}.
\]
Assign a monomial containing $t$ of the $y$'s exterior degree $|S|+2t$.
On each summand set
\[
 \mathscr L=\sum_{a\notin S}y_a\partial_{x_a},\qquad
 \mathscr D=\sum_{a\notin S}x_a\partial_{y_a}.
\]

\paragraph{Statement IF-01.}
There is a graded integral isometry $\Phi_n:\mathcal F_n\longrightarrow V_n$
intertwining $(\mathscr L,\mathscr D)$ with $(L,D)$ and all their divided
powers. It is given, with factors in increasing pair order, by
\[
 \Phi_n\left(\prod_{a\notin S\cup T}x_a\prod_{a\in T}y_a\right)
 =\prod_{a=1}^n u_a,\qquad
 u_a=\begin{cases}\epsilon_a&a\in S,\\p_a&a\in T,\\1&\text{otherwise}.
 \end{cases}
\]
In particular this is a specified lattice-preserving Fock realization of
Lefschetz divided powers, not merely an equality of rational characters.

\begin{proof}
A standard exterior basis vector chooses, in each symplectic pair, one of
$1,e_a,f_a,p_a$. The displayed assignment is therefore a bijection of
integral bases and preserves their declared pairings and degrees.
The operator $y_a\partial_{x_a}$ replaces the unique factor $x_a$ by
$y_a$, if present, and otherwise kills the monomial. Under $\Phi_n$ this
inserts $p_a$ in an empty pair and vanishes on a singleton or complete
pair. Because $p_a$ has even degree, moving it to its prescribed pair
position introduces no sign. Similarly
$\iota_{f_a^\vee}\iota_{e_a^\vee}$ removes a complete pair with sign $+1$,
and kills the other three states. This proves the two intertwining
identities. The commuting operators for distinct pairs square to zero;
hence
\[
 L^r/r!=\sum_{|R|=r}\prod_{a\in R}(p_a\wedge-)
\]
is integral, and the corresponding expression holds for $D$ and for the
two polynomial operators. Equivalently the coefficient of $z^r$ in the
simultaneous row substitution $x_a\mapsto x_a+zy_a$ is
$\mathscr L^r/r!$. Thus the isometry also respects the actual coefficient
normalization of divided row transvections.
\end{proof}

The exact source comparison is Alp\"oge,
\texttt{Formalism/GradedWeylDivisibility/StatementA.lean}, declaration
\texttt{rowDiv}: it is coefficient extraction under a row transvection.
The present module is its explicitly specified multiaffine specialization.
It does not identify a global theta-lift integral with that transvection.
On the exterior side this is the model in cote 16, scan 014, paragraph 7;
its contraction must remain distinguished from the inverse identities on
scan 016, as HODGE\_SIGN SL-04 already demonstrates.

\subsection{IF-02: the divided-power order is saturated for every rank}
This is the normalization-specific type-$A_1$ coefficient-coalgebra
statement for the actual exterior lattice. We claim no novelty for
the general integral Schur-algebra framework.
Write
\[
 \mathcal O_n=\mathbb Z[L,D],\qquad
 \mathcal B_n=\mathbb Z[L^{(r)},D^{(r)}:r\geq0],\qquad
 \mathcal A_n^{\rm ext}=\mathbb Q[L,D]\cap\operatorname{End}_{\mathbb Z}(V_n),
\]
where $L^{(r)}=L^r/r!$ and similarly for $D$; all powers with $r>n$
vanish. The intersection is taken inside
$\operatorname{End}_{\mathbb Q}(V_n\otimes\mathbb Q)$.

\paragraph{Statement IF-02.}
For every $n\geq0$,
\[
 \mathcal B_n=\mathcal A_n^{\rm ext},\qquad
 \operatorname{rank}_{\mathbb Z}\mathcal A_n^{\rm ext}
 =\sum_{v=0}^n(v+1)^2=\frac{(n+1)(n+2)(2n+3)}6.
\]
All exterior degree projectors belong to $\mathcal B_n$. For $n\geq1$
the additive index
$[\operatorname{End}_{\mathbb Z}(V_n):\mathbb Z[L,D]]$ is infinite.
Consequently the finite order-index question in the intake has to be
posed in $\mathcal A_n^{\rm ext}$, not in the whole endomorphism ring.

\begin{proof}
We give a coefficient-lattice proof, including the integral saturation
step. Put $U=\mathbb Z^2$ with its standard $\mathrm{SL}_{2,\mathbb Z}$
action. IF-01 identifies $V_n$, for this action, with a direct sum of
copies of $U^{\otimes v}$, $0\leq v\leq n$: the multiplicity is
$\binom n{n-v}2^{n-v}$. Every $v$ occurs. Repeated identical blocks do
not change either the image order or its saturation, so we may use
$W_n=\bigoplus_{v=0}^n U^{\otimes v}$.

Let
\[
 R=\mathbb Z[a,b,c,d]/(ad-bc-1),\qquad
 R_n=\langle\text{monomials of total degree at most }n\rangle_{\mathbb Z}.
\]
The matrix coefficients of $W_n$ span exactly $R_n$: the coefficients
of $U^{\otimes v}$ are all degree-$v$ monomials in the four matrix
entries. The relation $ad=bc+1$ reduces every element to a linear
combination of monomials not containing both $a$ and $d$. These normal
monomials are independent, as will also follow from the triangular
calculation below. The number of normal monomials of degree $v$ is
\[
 \binom{v+3}{3}-\binom{v+1}{3}=(v+1)^2.
\]
The coproduct of $R$ carries $R_n$ into $R_n\otimes R_n$. Therefore
$S_n=\operatorname{Hom}_{\mathbb Z}(R_n,\mathbb Z)$ is an algebra, acting
faithfully on $W_n$ by evaluating its matrix coefficients. Faithfulness
uses that those coefficients span $R_n$.

We show that the divided root operators span $S_n$ integrally. On the
big cell write
\[
 g=\begin{pmatrix}1&0\\x&1\end{pmatrix}
   \begin{pmatrix}\tau&0\\0&\tau^{-1}\end{pmatrix}
   \begin{pmatrix}1&y\\0&1\end{pmatrix}
 =\begin{pmatrix}\tau&\tau y\\\tau x&\tau xy+\tau^{-1}\end{pmatrix}.
\]
The normal monomials have the following expansions:
\[
 a^r b^s c^t\mapsto x^t y^s\tau^{r+s+t}\quad(r\geq0),
\]
\[
 b^s c^t d^u\mapsto
 \sum_{j=0}^u\binom uj x^{t+j}y^{s+j}\tau^{s+t-u+2j}
 \quad(u\geq1).
\]
For the first family, the exponent of $\tau$ minus the sum of the
exponents of $x,y$ is $r\geq0$. For the second it is the fixed negative
integer $-u$. In the second family the $j=0$ term has coefficient one;
all other terms increase $s+t$ by $2j$. Thus the coefficient submatrix
on the $j=0$ monomials, ordered by $s+t$, is triangular with diagonal
one, separately for each $u$. The first family contributes an identity
matrix. This proves independence of the normal monomials and shows that
$R_n$ is a primitive sublattice of the Laurent-polynomial coefficient
lattice. More precisely, restriction of coefficient-extraction
functionals spans the full dual $S_n$. Every appearing power of
$\tau$ lies between $-n$ and $n$, and the powers of $x,y$ lie between
zero and $n$.

Let $E=D$, $F=L$, and let $H$ act on $U$ with weights $+1,-1$.
Coefficient extraction in $x^r y^s z^k$ after setting $\tau=1+z$
acts as
\[
 F^{(r)}\binom Hk E^{(s)}.
\]
For exponents $q=-n,\ldots,n$, the functions $q\mapsto\binom qk$,
$0\leq k\leq2n$, form an integral basis of all functions on this finite
set. Indeed, $\binom{q+n}{k}$ gives a lower triangular matrix with
diagonal one, and Vandermonde's identity changes it to $\binom qk$
by an integral triangular matrix with diagonal one. Hence extraction of
an individual Laurent coefficient $\tau^q$ is an integral combination
of these coefficient functionals.

For completeness, the toral operators $\binom Hk$ lie in the algebra
generated by the divided root operators. The matrix identity
\[
 \begin{pmatrix}1&u\\0&1\end{pmatrix}
 \begin{pmatrix}1&0\\v&1\end{pmatrix}
 =\begin{pmatrix}1&0\\v/(1+uv)&1\end{pmatrix}
  \begin{pmatrix}1+uv&0\\0&(1+uv)^{-1}\end{pmatrix}
  \begin{pmatrix}1&u/(1+uv)\\0&1\end{pmatrix}
\]
expanded as a formal power series at $u=v=0$ gives
\[
 E^{(r)}F^{(s)}=
 \sum_{j=0}^{\min(r,s)}F^{(s-j)}
 \binom{H-r-s+2j}{j}E^{(r-j)}. \tag{IF-02a}
\]
To check the coefficient, choose powers $s-j$ and $r-j$ from the
left and right unipotent factors; their two denominators leave
$(1+uv)^{H-r-s+2j}$, and its $j$th coefficient is the displayed
binomial. For $r=s=k$, the final term is $\binom Hk$. Every earlier
term involves a shifted binomial of degree $j<k$, an integral
combination of $\binom Hh$, $h\leq j$, by Vandermonde's identity
(including its integral negative-shift form). Induction on $k$ proves
the assertion. Thus all coefficient functionals just constructed belong
to $\mathcal B_n$, proving $\mathcal B_n=S_n$ as embedded orders.

Over $\mathbb Q$, ordinary and divided powers generate the same algebra,
so $\mathbb Q[L,D]=S_n\otimes\mathbb Q$. A rational functional on
$R_n$ acts integrally on $W_n$ if and only if it is integer-valued on
every matrix coefficient, equivalently on $R_n$. This proves
$\mathcal A_n^{\rm ext}=S_n=\mathcal B_n$. A degree projector is the
indicator of one weight of $H$, since exterior degree $j$ has weight
$n-j$. The same integral binomial interpolation puts it in
$\mathcal B_n$. Finally $\operatorname{End}_{\mathbb Z}(V_n)$ has rank
$(4^n)^2=16^n$, strictly larger than $\sum_{v=0}^n(v+1)^2$ for
$n\geq1$. The rank comparison proves that its index over $\mathcal O_n$
is infinite.
\end{proof}

This proof also identifies a precise integral object in Grothendieck's
language: the divided-power image is the dual of the finite coefficient
coalgebra $R_n$ of the split group scheme $\mathrm{SL}_{2,\mathbb Z}$.
The exterior module is $\bigoplus_i\lambda^i(M_n)$ in the representation
$\lambda$-ring. No rational decomposition into primitive strings was
used to choose an integral lattice.

\subsection{IF-03: the actual exterior order indices through rank four}
Let $\mathcal O_n^\pi=\mathbb Z[L,D,\pi_0,\ldots,\pi_{2n}]$.
The exact indices in the canonical saturation are
\[
\begin{array}{c|r|c|c|c}
 n&\operatorname{rank}\mathcal A_n^{\rm ext}
 &[\mathcal A_n^{\rm ext}:\mathcal O_n]
 &[\mathcal A_n^{\rm ext}:\mathcal O_n^\pi]
 &[\mathcal A_n^{\rm ext}:\mathcal B_n]\\\hline
 1&5&1&1&1\\
 2&14&2^7&2^2&1\\
 3&30&2^{27}3^4&2^{10}3^2&1\\
 4&55&2^{81}3^{33}&2^{32}3^{10}&1
\end{array} \tag{IF-03a}
\]
The nonzero embedded Smith factors, including units, are
\[
\begin{array}{c|l|l}
 n&\mathcal O_n&\mathcal O_n^\pi\\\hline
 1&1^5&1^5\\
 2&1^8,2^5,4&1^{12},2^2\\
 3&1^8,2^{17},4,12^4&1^{20},2^8,6^2\\
 4&1^8,2^{19},6^5,12^{14},24^4,72^4,288
   &1^{28},2^{17},6^7,12,24^2.
\end{array} \tag{IF-03b}
\]
Here exponents on the Smith factors mean multiplicities. These are
indices on the actual exterior lattice. They are not the indices
$1,2^5,2^{13}3^{10}$ on the different ordinary-power string lattice
of HODGE\_SIGN SL-09.

\begin{proof}[Exact computation certificate, with an elementary rank-two check]
We specify both the finite algorithm and its certificate. In IF-01,
retain one copy of each singleton block, that is, the direct sum of the
Boolean modules on subsets of $\{1,\ldots,v\}$, $0\leq v\leq n$.
There are $\sum_{v=0}^n4^v$ matrix coordinates. On the $v$th block,
\[
 L_{T,S}=\begin{cases}1&S\subset T,\ |T\setminus S|=1,\\0&\text{otherwise},\end{cases}
 \qquad D=L^{\mathsf t},\qquad
 (\pi_j)_{S,S}=\boldsymbol1_{n-v+2|S|=j}.
\]
Repeated singleton blocks have identical integer matrices and hence do
not affect integral saturation. Starting from the column vector of the
identity, repeatedly take the column Hermite normal form of
\[
 [B,\ R_LB,\ R_DB]
 \quad\text{or}\quad
 [B,\ R_LB,\ R_DB,\ R_{\pi_0}B,\ldots,R_{\pi_{2n}}B],
\]
where $R_g$ means right multiplication by the specified integral matrix.
By induction this is exactly the lattice spanned by words in the
generators of bounded length. At equality of two successive HNF bases
it is closed under every generator, contains the identity, and is
therefore exactly the generated order. Equality is equality of integer
matrices, not just equality of ranks. Applying Smith reduction to this
embedded lattice gives its index in its saturation as the product of the
nonzero Smith factors. IF-02 identifies that saturation independently.

The complete bases, the stabilization histories, and their Smith factors
are the files \texttt{order\_basis\_nN\_KIND.json}; the algorithm is
\texttt{exterior\_order\_indices.py}. Each of the sixteen computed orders
is checked for closure under every generator. The independent program
\texttt{verify\_exterior\_orders.py} constructs the full signed exterior
matrices directly, verifies the generator normalization, and checks the
nontrivial local determinantal valuations by $p$-local elimination at
$2,3,5$, without using the Smith routine. Its output is
\texttt{independent\_exterior\_verification.json}. A separate check
by the root agent, \texttt{work/root/verify\_left\_closure.py},
starts from independently constructed full signed exterior matrices,
generates left words, and compares all sixteen complete HNF lattices
for $n=1,2,3,4$ with the right-word certificates. It passed; its
output is \texttt{work/root/left\_closure\_verification.json}. The
two closure computations share SymPy HNF infrastructure. The factors in
(IF-03b) are the exact integer outputs of the stated terminating
algorithm; multiplication gives (IF-03a). This is a computation
certificate, not a Lean formalization.

For an elementary check on the normalization-sensitive case $n=2$, take
the three-dimensional string with basis
$v_0=1$, $v_1=p_1+p_2$, $v_2=p_1p_2$, and the trivial vector
$a=p_1-p_2$. On this string
$Lv_0=v_1$, $Lv_1=2v_2$, $Dv_1=2v_0$, $Dv_2=v_1$.
Write $A=(a_{rs})_{0\leq r,s\leq2}$ for its rational operator and
$z$ for the operator on $a$. The independent trivial singleton
block forces $z\in\mathbb Z$. Preservation of the lattice with
$p_1=(v_1+a)/2$, $p_2=(v_1-a)/2$ is exactly
\[
 a_{01},a_{21}\in2\mathbb Z,\qquad
 a_{11}\equiv z\pmod2,
\]
with all other matrix entries integral. This is the even part of
$\mathcal A_2^{\rm ext}$; the odd part is a full $2$-by-$2$ matrix ring.
The degree projectors and $L,D$ generate this even order with the two
additional conditions $a_{02},a_{20}\in2\mathbb Z$. Indeed the endpoint
projectors generate
$E_{10},2E_{21},E_{12},2E_{01},2E_{20},2E_{02}$,
and their products give $2E_{11}$; the middle projector supplies
$(E_{11},1)$. These are precisely a basis subject to the displayed
conditions and the two additional divisibilities. Thus the index with
degree projectors is $2^2$, agreeing with the certificate and showing
explicitly why the string-model value $2^5$ does not apply.
\end{proof}

\subsection{IF-04: the full Weyl lattice and its dual in an exterior model}
The multiaffine construction can be extended to the actual polynomial
Weyl lattices in the source. This also locates their denominators.

Let $U=\mathbb Z^r$ and let $\lambda=(\lambda_1\geq\cdots\geq
\lambda_r\geq0)$ be a partition. Let $h_1,\ldots,h_q$ be its column
heights, where $q=\lambda_1$. Set
\[
 T^{\rm amb}_\lambda=\bigotimes_{j=1}^q\bigwedge^{h_j}U,
 \qquad t_\lambda=\bigotimes_{j=1}^q(e_1\wedge\cdots\wedge e_{h_j}).
\]
Give this tensor lattice the orthonormal exterior-basis pairing. In the
polynomial ring on an $r$-by-$r$ matrix let $\Delta_h$ be its upper-left
$h$-minor and put $p_\lambda=\prod_j\Delta_{h_j}$.
Let $T_\lambda$ be the lattice generated from $t_\lambda$ by all divided
row-root operators and let $\Delta_\lambda$ be the lattice generated
from $p_\lambda$ by those operators. The latter is exactly
\texttt{weyl} in \texttt{GradedWeylDivisibility/StatementA.lean},
with \texttt{hw} as its initial vector. A determinant twist is removed
as in \texttt{LocalRingCardinality/DataCoeff.lean}, declarations
\texttt{DomWeight.detShift} and \texttt{DomWeight.natW}.

\paragraph{Statement IF-04.}
Multiplying the initial minors gives an integral isomorphism
\[
 \pi:T_\lambda\xrightarrow{\ \sim\ }\Delta_\lambda.
\]
It is an isometry when the polynomial Fock pairing is divided by
$\langle p_\lambda,p_\lambda\rangle_{\rm Fock}$, precisely the source's
\texttt{normFockQ}. Fix a simple row root $i,i+1$ and put
$d=\lambda_i-\lambda_{i+1}$. Its cyclic string has basis
\[
 q_s=F_i^{(s)}t_\lambda\quad(0\leq s\leq d),
\]
with
\[
 \langle q_s,q_t\rangle=\boldsymbol1_{s=t}\binom ds,
 \quad F_i^{(k)}q_s=\binom{s+k}{k}q_{s+k},
 \quad E_i^{(k)}q_s=\binom{d-s+k}{k}q_{s-k},                 \tag{IF-04a}
\]
where out-of-range vectors are zero. The weight-line lattice dual to
$\mathbb Z q_s$ is $\mathbb Z q_s/\binom ds$. Moreover, this entire
root action has an explicit integral realization by ordinary Lefschetz
operators on a symplectic exterior lattice; its Weyl string is primitive
in that exterior lattice. The dual string generally is not contained in
that same exterior lattice.

\begin{proof}
Map $e_{i_1}\wedge\cdots\wedge e_{i_h}$ in a column factor to the minor
with rows $i_1,\ldots,i_h$ and the first $h$ columns, and multiply these
polynomials over all column factors. This defines an integral linear map
$\pi:T^{\rm amb}_\lambda\to\mathbb Z[X_{ab}]$, equivariant for row
substitutions, sending $t_\lambda$ to $p_\lambda$.

We prove injectivity on the specified cyclic lattice without presuming
an integral decomposition into irreducibles. For each initial set of
$k$ rows, a wedge of degree $h$ has at most $\min(k,h)$ entries in those
rows. The row weight of $t_\lambda$ attains the sum of these bounds for
every $k$. Equality for every $k$ forces each wedge factor to be its
initial wedge. Thus its weight line in $T^{\rm amb}_\lambda\otimes\mathbb Q$
is one-dimensional. Ordinary opposite row-root operators are adjoint
for the tensor exterior pairing and also for the factorial polynomial
Fock pairing. Since $\pi$ intertwines both roots, its rational adjoint
$\pi^*$ does also, and $A=\pi^*\pi$ commutes with every root and Cartan
operator. The one-dimensional weight line gives
$At_\lambda=c t_\lambda$, with
\[
 c=\langle p_\lambda,p_\lambda\rangle_{\rm Fock}>0,
\]
because $t_\lambda$ has norm one and $p_\lambda\ne0$.
Commuting with every divided root word shows that $A=c$ on
$T_\lambda\otimes\mathbb Q$. Thus $\pi$ is injective there, and it
multiplies the pairing by $c$. Its integral image is exactly
$\Delta_\lambda$ by the defining row-word generators. This proves the
integral isomorphism and the normalized isometry.

Exactly $d$ column factors have height $i$. The root $F_i$ can change
$e_i$ to $e_{i+1}$ only in these factors; it is zero on the initial
wedges of all other heights. Consequently $q_s$ is the sum, with
coefficient one, over the $s$-subsets of these $d$ columns. Its summands
are distinct orthonormal basis vectors. This proves the norm formula.
Counting the predecessors or successors of a fixed subset proves the
two binomial operator formulas. The sums for different $s$ have
disjoint supports and each has a coefficient one, so their span is a
primitive sublattice of the ambient tensor lattice.

We also check that this is the complete relevant weight-line lattice of
$T_\lambda$. Over $\mathbb Q$, root words can be reordered with negative
roots first, Cartan operators next, and positive roots last: exchanging
two factors uses their commutator, reducing the word length, so induction
performs the reordering. Positive roots kill $t_\lambda$. To reach the
weight $\lambda-s\alpha_i$, every remaining negative root must be
$\alpha_i$, because each positive root is a nonnegative sum of simple
roots. Hence that rational weight space is the one-dimensional span
of $q_s$. Primitivity in the ambient tensor lattice gives its integral
intersection $\mathbb Z q_s$. Different weights are orthogonal. The
normalized dual lattice therefore has intersection
$\mathbb Z q_s/\binom ds$ with this rational line. This is the Weyl/dual
Weyl pairing used in \texttt{StatementE.lean}, declarations
\texttt{NablaEqDualWeyl} and \texttt{WeylEqDualNabla}, and in
\texttt{DataCoeff.lean}, declaration \texttt{FactDelDualNabla}.

Here is the promised exterior construction. Embed
$T^{\rm amb}_\lambda$ into
$\bigwedge^*(\bigoplus_{j=1}^q U_j)$ by putting the $j$th wedge in the
$j$th summand. In each column order the root coordinates $e_i,e_{i+1}$
first. On their four exterior states apply the integral basis
permutation
\[
 e_i\mapsto1,\qquad e_{i+1}\mapsto e_i\wedge e_{i+1},\qquad
 1\mapsto e_i,\qquad e_i\wedge e_{i+1}\mapsto e_{i+1}.
\]
It carries the row operator $e_{i+1}\wedge\iota_{e_i^\vee}$ to wedge
multiplication by the complete root pair, and the opposite root to its
contraction. Thus the sum over columns becomes a Lefschetz $L$, and the
opposite sum becomes $D$. The remaining $(r-2)q$ exterior coordinates
are spectators. For a nondegenerate symplectic ambient space, replace
each spectator's two states by the two singleton states of a fresh
symplectic pair; both are killed by its $L,D$. This is an integral
isometric embedding into $V_{(r-1)q}$, and the full symplectic $L,D$
restrict to the desired root operators. The displayed $q_s$ still have
disjoint supports with coefficients one. Therefore they are primitive
in this actual exterior lattice. If $\binom ds>1$, the vector
$q_s/\binom ds$ is not integral there.
\end{proof}

In particular factorial Fock positivity by itself does not choose the
Weyl lattice instead of its dual. For one binary column of degree $d$,
the polynomial lattice has monomial basis $x^{d-s}y^s$ whereas the
row-divided orbit of $x^d$ has basis
$\binom ds x^{d-s}y^s$. Both inherit the same positive factorial
pairing. IF-04 identifies the latter with orbit sums on the exterior
side; the former becomes its normalized dual. Thus the sentence in
\texttt{INTAKE\_HODGE\_SIGN\_20261121.md}, section 5, identifying a
factorial pairing alone with the binomial lattice needs this refinement.

The cyclic lattice is weight-graded over $\mathbb Z$: every divided-root
word on the highest tensor is weight-homogeneous. Thus orthogonality
computes the intersection with the full dual lattice, not just a dual
inside an isolated rational string. The source also permits
\texttt{rowDiv a a k}; on a homogeneous row weight $w_a$ this is
multiplication by $\binom{w_a}{k}$. These diagonal coefficient
operators therefore add nothing to the cyclic integral lattice.

\subsection{IF-05: the auxiliary row-factor test at two lifts}
The actual column-kernel calculation, with determinant shifts and
finite theta jets, is IF-08--IF-09 below. This subsection tests the
distinct auxiliary compact row factors used in the same source.
We use the actual numerical datum \texttt{numEx}, not an invented
rank-two theta tower. Its declarations are
\texttt{LocalToGlobal/Divisibility/NumbersLemmas.lean}, namespace
\texttt{Formalism.L2.DIV}, declaration \texttt{numEx};
\texttt{LocalRingCardinality/SpanPowerMembership/ExampleRanksRankComplement.lean},
\texttt{R0}; and
\texttt{ExampleWeightMiddleEntries.lean}, \texttt{hcV,hcB,hcE}.
The data are $m=2,k=8,t=15,nV=24$, $e=(100,118)$, $b=(91,109)$,
$v=(24,20,16,12,8,4,0)$, and the definite tower ranks are $8,24,40$.
The auxiliary signatures \texttt{sigW} are $(8,1)$ and $(9,0)$.
The row functions \texttt{aPlus,aMinus} are defined in
\texttt{Numbers.lean}; their actual values are
\[
\begin{array}{c|l|l}
 j&a_j^+&a_j^-\\\hline
 1&(109,106,103,100,97,94,91,1)&(1)\\
 2&(126,123,120,117,114,111,108,18,1)&()
\end{array}.
\]
The adjacent positive-row differences are
$(3,3,3,3,3,3,90)$ and $(3,3,3,3,3,3,90,17)$.
Constant determinant shifts do not change them.

\paragraph{Statement IF-05.}
Apply IF-04 to the polynomial Weyl factors prescribed by these row
weights. On an actual simple-root string of length $d+1$, two successive
ordinary lowering operators have Smith invariant $2$ from the top line
to the second descendant in the Weyl lattice, and $d(d-1)$ in its
normalized dual. The second divided operator has invariants $1$ and
$\binom d2$, respectively. In particular the actual available values are
\[
\begin{array}{c|c|c|c}
 d&\text{two ordinary steps: Weyl, dual}
   &\text{second divided step: Weyl, dual}&\binom d2\\\hline
 3&(2,6)&(1,3)&3\\
 90&(2,8010)&(1,4005)&3^2\cdot5\cdot89\\
 17&(2,272)&(1,136)&2^3\cdot17
\end{array}.                                                   \tag{IF-05a}
\]
The isometry of IF-04 intertwines the Weyl operators integrally with
Lefschetz operators on its explicitly constructed exterior sublattice.
It cannot simultaneously carry the normalized dual lattice into that
same exterior lattice. More generally, for $d>1$ the normalized dual
root string admits no isometric integral injection into any lattice
having an orthonormal integral basis: it contains a nonzero vector of
squared norm $1/d<1$.

\begin{proof}
In the Weyl basis $q_s$ of IF-04, the first two ordinary arrows have
coefficients $1,2$; their product is $2$. In the normalized dual basis
$r_s=q_s/\binom ds$, the first two coefficients are $d,d-1$. Division
of the composite by $2!$ gives the stated divided invariants. These are
maps between free rank-one lattices, so the absolute coefficient is
exactly the Smith invariant. Substitution of the three adjacent gaps
gives the table. The norm calculation gives
$\langle r_1,r_1\rangle=1/d$. A nonzero integral vector in an
orthonormal integral lattice has squared norm a positive integer. This
proves the obstruction. It also shows concretely why replacing a Weyl
lattice by its Fock dual changes an index despite preserving the
rational representation and the root operator.
\end{proof}

The fixed row-weight test therefore has both a positive and a negative
outcome: there is an explicitly constructed integral normalized Weyl
intertwiner, and there are explicitly computed nontrivial denominators
for the simultaneously required dual. The normalizations have not been
suppressed. This advances HODGE\_SIGN U-04 from unspecified Fock data to
an actual lattice realization and an exact denominator calculation.

For comparison with the arithmetic entries themselves, HODGE\_SIGN
SL-06's exact substitution in \texttt{Numbers.lean}, declarations
\texttt{Dplus,Dminus,FTypeG,FRamD,FInert,FjP}, gives
\[
 (D_1^+,D_2^+)=(5608,6704),\quad(D_1^-,D_2^-)=(8,0),
\]
\[
 (F_{G,1},F_{G,2})=(-2808,-3352),\quad
 (F_{{\rm ram},1},F_{{\rm ram},2})=(-1400,-1652),\quad
 F_{{\rm inert},j}(p)=\frac{56e_j}{p^2-1}.
\]
These particular substitutions are rechecked here; the orbit optimization
underlying the last formula is cited as SL-06, not proved again.
At $p=3$ the two arithmetic contributions are $700,826$, while at
$p=89$ both \texttt{FjP} values are zero: $89$ exceeds the datum's cutoff
$X''=8\exp(2^{1/300})<22$ and is neither a distinguished prime nor the
ramified prime $7$. Nevertheless the fixed Weyl/dual comparison for its
$d=90$ row string has a factor $89$. Thus an asserted equality of the
raw step-10 $F$ contribution with this Fock lattice index is false on
the actual $m=2$ data. The full step-10 theorem, which additionally
contains gains, errors and stabilizer denominators, is not refuted.
Here $89$ differs from $3697,7393$ (the distinguished primes),
$7$ (the ramified $d$ prime), and also $2,11$; the above-cutoff
lemma only requires the first three exclusions. The exact statement is
\texttt{NumbersLemmas.lean}, \texttt{FjP\_above}, lines 191--194.
Moreover \texttt{numEx} itself is not a full \texttt{NumFacts}
witness. The change to the valid \texttt{numBig} instance and its
column weights is recorded explicitly in IF-09; its larger cutoff
is not being substituted into this $89$ comparison.
A reproducible rational-arithmetic certificate is
\texttt{step10\_fock\_test.py} and its tables.

\paragraph{Full-module rank check.}
The positive compact factors determined by the two row shapes, after
any constant determinant twist, have Weyl dimensions
\[
 105068200743744433029120,
 \qquad17540857265584852885502309695488000,
\]
respectively. Both have odd prime divisors; for example the first has
factor $3$ and the second has factor $3^4$. Consequently neither factor
can be isomorphic to a full $V_n$, whose rank is $4^n$. These numbers
come from the Weyl dimension product
$\prod_{a<b}(\lambda_a-\lambda_b+b-a)/(b-a)$ (Etingof et al.,
\emph{Introduction to representation theory}, Theorem 4.63,
\url{https://math.mit.edu/~etingof/replect.pdf}), evaluated in exact
rational arithmetic in the same certificate. This rank test concerns
these finite compact factors. It does not assert that the rank of a
global automorphic lattice is the rank of this one factor. IF-04 gives
a suitable sublattice realization rather than a full-exterior
identification.

\subsection{IF-06: the full oscillator cannot be the finite exterior module}
The preceding restricted construction does not justify replacing the
whole oscillator representation by its multiaffine part. The actual
source makes the distinction visible.

In \texttt{LocalRingCardinality/DataTheta.lean}, the definitions
\texttt{lowB,raiB,omE,fockAct1} and theorem
\texttt{fockAct1\_Eab} imply the following. At a signature $(p,q)$ with
$p,q>0$, against a positive definite partner of rank $R>0$, choose one
positive and one negative row. The noncompact raising root on the
polynomial Fock space is multiplication by
\[
 Q=\sum_{c=1}^R x_c y_c,
\]
and the opposite root is $-C$, where
$C=\sum_c\partial_{x_c}\partial_{y_c}$. On these variables
\[
 [C,Q]=H=R+\operatorname{Euler},\qquad [H,Q]=2Q,\qquad[H,C]=-2C.
\]
This applies in particular to the actual auxiliary signature
$\texttt{sigW}(1)=(8,1)$ of the $m=2$ datum.

\paragraph{Statement IF-06.}
Over any characteristic-zero field, there is no nonzero intertwiner
from a finite-dimensional Lefschetz module to this polynomial Fock
space taking its raising operator to $Q$. If both roots and their
commutator are intertwined, there is also no nonzero intertwiner in
the opposite direction, even after replacing contraction by its
negative to match the $\mathfrak{sl}_2$ signs.

\begin{proof}
Multiplication by the nonzero polynomial $Q$ is injective, and every
power of it is injective. On a finite Lefschetz module $L^{n+1}=0$.
Thus an intertwiner $T$ into the polynomial space satisfies
$Q^{n+1}T=TL^{n+1}=0$ and is zero.

For the converse, first decompose the polynomial space into strings
$Q^a h$ with $Ch=0$. Here is a direct existence proof of that familiar
harmonic decomposition. If $h$ is homogeneous of degree $d$ and
$Ch=0$, induction from the commutators gives
\[
 C Q^a h=a(R+d+a-1)Q^{a-1}h.                       \tag{IF-06a}
\]
For a homogeneous polynomial of degree $j$, apply induction to its
image under $C$, which has degree $j-2$. Each string term $Q^a h$
in that image has a preimage
$Q^{a+1}h/((a+1)(R+\deg h+a))$; the denominator is a nonzero positive
integer. Subtracting these preimages leaves a harmonic polynomial.
This proves that the stated strings span, over the characteristic-zero
field. Other polynomial variables may be treated as fixed spectator
coefficients.

Let $T$ go from the oscillator to a finite Lefschetz module, with
$TQ=LT$ and $T(-C)=DT$. For a harmonic homogeneous $h$, the vector
$w=Th$ is killed by $D$ and has eigenvalue $-(R+\deg h)$ for
$H_{\rm Lef}=[D,L]$. In a finite Lefschetz module a nonzero vector
$w$ killed by $D$ cannot have negative $H_{\rm Lef}$ weight. To see
this without a decomposition theorem, let $L^a w\ne0$ be its last
nonzero translate and write $H_{\rm Lef}w=\delta w$. The commutator
formula gives
$0=DL^{a+1}w=(a+1)(\delta-a)L^a w$, so $\delta=a\geq0$.
Here the required weight is $-(R+\deg h)<0$, a contradiction.
Therefore $Th=0$, and $TQ^a h=L^aTh=0$ on every spanning string.
\end{proof}

This rules out the full-oscillator shortcut. It does not rule out
finite compact pieces or the explicitly constructed integral
intertwiner IF-04. In particular it is not a nonexistence theorem for
every possible geometric realization of the global definite tower
$8\to24\to40$.

\subsection{IF-07: what an exact period valuation can depend on}
For fixed integral linear maps $K_0,K_1$ and an integral terminal
functional $\ell$, put $P(v)=\ell(K_1K_0v)$. If $P(v)\ne0$, then
for every rational prime $p$ and $s\geq0$,
\[
 v_p(P(p^s v))=s+v_p(P(v)).
\]
This follows immediately from linearity, and the same identity holds
for linear theta lifts and their terminal periods when the allowed
input vector is scaled. Thus lift count alone does not determine a
nonzero period's exact valuation. With a fixed input vector this
argument says nothing about the valuation; it is not used to dismiss
that stronger problem.

There is already a fully constructed nonzero local comparison here:
for the actual root strings of IF-05 take $v=q_0$ and pair the result
with its normalized Fock dual $r_2$. The ordinary two-step period is
$\langle r_2,F_i^2q_0\rangle=2$, and the divided two-step period is
$\langle r_2,F_i^{(2)}q_0\rangle=1$. All their valuations are exact.
These are periods of the specified row-root maps, not evaluations of
the global theta integral. Identifying the latter with these maps
would require a further equality of actual kernels and their
stabilizer normalizations. \texttt{Inputs.lean}, structure
\texttt{ColChain}, supplies the rectangular global matrices and a
factorization, but no such equality is an axiom of that structure.

\paragraph{Algebraic-geometric placement and precise remaining assertion.}
IF-01--IF-05 live over $\operatorname{Spec}\mathbb Z$ and commute with
base change. For a principally polarized complex abelian variety,
$H^*(A,\mathbb Z)=\bigwedge^*H^1(A,\mathbb Z)$ realizes their exterior
operator calculations, with $L$ cup product by the polarization.
The same matrices act on prime-to-characteristic integral
$\ell$-adic cohomology of a principally polarized abelian scheme.
Passing from degree $i$ to $i+2$ uses a Tate twist; after the twists in
HODGE\_SIGN SL-07, the correspondence commutes with Frobenius. The
calculation does not make an integral inverse Lefschetz correspondence
algebraic, and it does not identify a Weyl lattice's Fock dual with an
integral cycle lattice. In $K$-theory the constructions use exterior
powers and tensor products; in the category of rational motives they
are operations on the appropriate realizations of $h^1(A)$, together
with the Lefschetz correspondence. No new identity of motivic
$L$-functions or of global theta periods follows from the lattice
isomorphisms alone.

The finite-kernel part of the predecessor's remaining assertion is
supplied next in IF-08--IF-09: each actual finite theta jet is expressed
by an explicit exterior-tensor matrix, including the dual lattices,
determinant shifts and stabilizer factors. What that does not assert
is that these varying tensor matrices are one common Lefschetz
$\mathfrak{sl}_2$ operator, or that an arbitrary global automorphic
period is the particular root period just computed. Those are distinct
operator and geometric assertions, not missing coefficient-model data.

% IF-08 and IF-09: actual coefficient kernels and actual m=2 columns.
% Included by contribution.tex. Full author proof reading recorded separately.
\subsection{IF-08: the actual finite theta jet has an exterior-tensor realization}
The finite coefficient map in the source is more general than a single
raising operator. It can nevertheless be realized explicitly and
integrally in exterior tensors, with exactly the same Smith indices.
Here the relevant source declarations are
\texttt{LocalToGlobal/KernelLattices/ThDistC.lean},
\texttt{subD,jetLin,jetMap,jetOp,JetInVV,derivAt0,jetC}
(lines 23--61), and
\texttt{LocalRingCardinality/DataKerI.lean},
\texttt{KM.Gxz,KM.applyAt0,KerEnt.kform} (lines 27--35 and 115--121).
The former is the coefficient map itself; the latter describes its
matrix entries tested against the dual product-minor lattice.

Fix a partition $\lambda$, with column heights $h_1,\ldots,h_q$,
and integers $a,b\geq\ell(\lambda)$. Put
\[
 A_\lambda(r)=\bigotimes_{c=1}^q\Lambda^{h_c}\mathbb Z^r,
 \qquad T_\lambda(r)\subset A_\lambda(r)
 \quad (r=a,b),
\]
where $T_\lambda(r)$ is the divided-root cyclic lattice of IF-04.
Write $\pi_r:A_\lambda(r)\longrightarrow\mathbb Z[X]$ for
multiplication of the corresponding initial-column minors. Use polynomial
matrices with $r$ rows and the common $\ell(\lambda)$ columns. This
makes rectangular substitution meaningful also when $b<a$; the
source \texttt{subD} sends unused source columns $s\geq b$ to zero,
and no such column occurs in these minors. Its cyclic
restriction is the isomorphism
$\Phi_r:T_\lambda(r)\xrightarrow\sim\Delta_\lambda(r)$ of IF-04.
The full tensor image is $\nabla_\lambda(r)$, by the supplied
\texttt{GradedWeylDivisibility/FinalC.lean}, theorem
\texttt{standardMonomialSpan}, whose exact statement is
\texttt{StatementA.lean}, \texttt{StandardMonomialSpan}, lines 221--226.
We use that supplied integral straightening result as background.

Let $Z=(z_{ui})$ be a $b$ by $a$ matrix of independent variables and set
\[
 A_\lambda(Z)=\bigotimes_c\Lambda^{h_c}Z
       =\sum_\beta Z^\beta A_\beta,
 \qquad A_\beta\in\operatorname{Hom}_{\mathbb Z}
       (A_\lambda(a),A_\lambda(b)).
\]
For a characteristic-zero coefficient field $K$ and a prescribed moment
function $J:\mathbb N^{ba}\to K$, define the concrete finite matrix
\[
 B_J=\sum_{|\beta|=|\lambda|}J(\beta)A_\beta.       \tag{IF-08a}
\]
Only finitely many entries of $J$ are used. If
$\Theta(Z)=\sum_\beta\theta_\beta Z^\beta$, the source convention is
\[
 J_\Theta(\beta)=\beta!\,\theta_\beta,
 \qquad \beta!=\prod_{u,i}\beta_{ui}!,              \tag{IF-08b}
\]
including the explicitly prescribed powers of the discriminant-period
generators in \texttt{jetC}. Thus $J$ consists of ordinary derivatives
at zero; it is not an unspecified divided-power normalization.

\paragraph{Statement IF-08.}
The lattices $T_\lambda(r)$ are saturated in $A_\lambda(r)$, and
\[
 \Phi_r:T_\lambda(r)\simeq\Delta_\lambda(r),\qquad
 \Phi_r:T_\lambda(r)^\#\simeq\nabla_\lambda(r)       \tag{IF-08c}
\]
are isometries for the tensor-exterior pairing on the left and the
normalized Fock pairing on the right. The matrix $B_J$ preserves the
rational cyclic spaces. If its moments lie in a subring $R\subset K$,
it carries $T_\lambda(a)\otimes R$ into $T_\lambda(b)\otimes R$.
On those spaces the actual polynomial substitution followed by moment
contraction satisfies
\[
 \Phi_b B_J=\operatorname{jet}_J\Phi_a.             \tag{IF-08d}
\]
This equation is an equality of explicit coefficient matrices. With
source and target lattices as displayed, it preserves every
Smith invariant over $\mathbb Z$ and every elementary-divisor exponent
over a discrete valuation ring. The transpose identity gives the same
statement for the normalized dual lattices.

For a determinant insertion of order $n\geq0$, multiplication by
$\det(X)^n$ is represented on the exterior side by insertion of $n$
copies of the oriented top wedge of $\mathbb Z^a$. Thus the complete
\texttt{jetMap J n}, including its determinant factor, is intertwined
by the same construction. Its rational retraction \texttt{retrVV}
has no effect under the supplied \texttt{JetInVV} hypothesis.

\begin{proof}
We first make the scalar normalization and saturation explicit. Let
\[
 c_\lambda=\langle p_\lambda,p_\lambda\rangle_{\rm Fock}
   =\sum_\alpha ([X^\alpha]p_\lambda)^2\alpha!>0.
                                                               \tag{IF-08e}
\]
This explicit finite sum is an integer independent of the ambient rank $r\geq\ell(\lambda)$.
IF-04 proves $\pi_r^*\pi_r=c_\lambda$ on $T_\lambda(r)_\mathbb Q$,
where the adjoint uses the unnormalized factorial Fock form and the
orthonormal exterior tensor form. The rational image of $\pi_r$ is
already the image of its cyclic subspace, by
\texttt{GradedWeylDivisibility/FinalB.lean}, theorem
\texttt{weylEqNablaQ} (the supplied \texttt{WeylEqNablaQ} statement).
Hence $T_\lambda(r)_\mathbb Q=(\ker\pi_r)^\perp$. On the whole tensor
space we therefore have
\[
 \pi_r^*\pi_r=c_\lambda P_r,\qquad
 P_r=\frac1{c_\lambda}\pi_r^*\pi_r,                 \tag{IF-08f}
\]
where $P_r$ is the orthogonal projection onto the cyclic rational
space. These formulas display, rather than discard, the scalar.

Suppose $t\in T_\lambda(r)_\mathbb Q\cap A_\lambda(r)$.
For every $v\in A_\lambda(r)$, (IF-08f) gives
\[
 \frac{\langle\pi_rt,\pi_rv\rangle_{\rm Fock}}{c_\lambda}
       =\langle t,v\rangle\in\mathbb Z.
\]
The full image $\pi_r(A_\lambda(r))$ is $\nabla_\lambda(r)$.
The supplied normalized duality theorem
\texttt{LocalToGlobal/Columns/CoefficientsModP.lean},
\texttt{weylEqDualNabla} (line 31),
therefore implies $\pi_rt\in\Delta_\lambda(r)$; IF-04 gives
$t\in T_\lambda(r)$. This proves saturation. The other normalized
duality, \texttt{nabla\_eq\_dual\_weyl\_holds} in
\texttt{Columns/CoefficientsModP/WeTransposeWordTransposeWord.lean}
(line 74), and the isometry of IF-04 yield
(IF-08c). Equivalently $P_r(A_\lambda(r))=T_\lambda(r)^\#$;
this also follows by extending an integral linear functional on the
primitive sublattice $T_\lambda(r)$ to the free ambient lattice.
The projection can have denominators, exactly those already present
in this specified dual lattice.

The exterior matrix $A_\lambda(Z)$ has entries
$\prod_c\det Z[I_c,J_c]$. Multiplication of minors gives
\[
 \pi_b A_\lambda(Z)=\rho_Z\pi_a,                  \tag{IF-08g}
\]
where $\rho_Z$ substitutes $X_{is}\mapsto\sum_u Z_{ui}X_{us}$.
For completeness, expand a determinant of these sums by multilinearity
in its rows. Terms with a repeated index vanish by alternation; grouping
all permutations of each increasing index set gives the determinant
$\det Z[I,J]$ times the corresponding minor of $X$. Multiplying these
identities over the columns gives (IF-08g), including all signs.

For a rational specialization of $Z$, (IF-08g) sends $\ker\pi_a$
into $\ker\pi_b$. Its transpose version sends $\ker\pi_b$ into
$\ker\pi_a$, since $A_\lambda(Z)^t=A_\lambda(Z^t)$.
Taking orthogonal complements shows
$A_\lambda(Z)T_\lambda(a)_\mathbb Q\subset
T_\lambda(b)_\mathbb Q$. This holds for every rational matrix $Z$.
Every coefficient $A_\beta$ consequently has the same containment:
apply any linear functional vanishing on the latter space and use
that a polynomial vanishing at all rational tuples is zero.
The coefficients $A_\beta$ are integer matrices; the saturation just
proved implies $A_\beta T_\lambda(a)\subset T_\lambda(b)$.
This proves the claimed integrality for arbitrary $R$ and $J$.

Apply the moment functional $Z^\beta\mapsto J(\beta)$ to (IF-08g).
It gives (IF-08d). For $J=J_\Theta$, direct differentiation of a
monomial proves (IF-08b); this is exactly the factorial in
\texttt{KM.applyAt0}. The test-minor version is
\[
 \operatorname{applyAt0}(\Theta,G_I(XZ))
   =\sum_J \Bigl(\sum_\beta\beta!\theta_\beta
                [Z^\beta]\prod_c\det Z[J_c,I_c]\Bigr)G_J(X),
                                                               \tag{IF-08h}
\]
with the direction of $Z$ transposed to match the source's test
convention. This is the same matrix, paired with the dual lattice.
Under (IF-08c), normalized Fock testing is ordinary tensor-exterior
testing. The $c_\lambda$ in (IF-08f) has canceled because the Fock
form is divided by that same $c_\lambda$; replacing normalized Fock
by unnormalized Fock would retain that factor.

Finally insert $n$ full-height columns in $A_\lambda(a)$. Their top
wedges have norm one and transform by $\det^n$, while every row root
annihilates them. Their cyclic sublattice is therefore exactly the
cyclic lattice of the shape $\lambda+n(1^a)$, and their minor product
is $\det(X)^n$. Thus determinant insertion also has no unspecified
scalar. In integral bases, (IF-08d) and its determinant variant are
conjugations by unimodular source and target matrices. Their images,
cokernels, saturated images and local lengths are correspondingly
isomorphic. This proves the assertions about Smith invariants,
including for rectangular maps and rank-deficient images, where the
finite index means the index in the saturation of the image.
\end{proof}

\paragraph{Assembly of the actual finite kernel.}
This is not merely a highest-weight-root comparison. For every
supplied theta jet $J_{j,c',c}$, (IF-08a) specifies its exterior-tensor
matrix. The complete block matrix retains the source's actual scalar
factors:
\[
 \widehat K_j(c',c)=\frac{\rho_j(c)\rho'_j(c')}{|\Gamma_{j,c}|}
        B_{J_{j,c',c}}\,D_{n_j}.                  \tag{IF-08i}
\]
Here $D_{n_j}$ inserts the determinant factors, and the target
one-dimensional determinant character is retained as a separate
rank-one lattice. The scalar entries and the jets include precisely
the period generators and frame changes specified in the source.
Formula (IF-08i) is the tensor version of
\texttt{SharpRestriction/GlobalSchurBounds/GL3Theta.lean},
\texttt{liftTie,lift\_formula} (lines 575--596), and
\texttt{LocalRingCardinality/DataKernels.lean},
\texttt{wt,sharp,divDe,Kmap,K} (lines 52--68).

The coefficient isomorphisms commute with the relevant group actions;
hence they transport invariant submodules, local lattice tests and
finite direct sums over lattice classes. For an actual supplied
finite global matrix they give an isomorphic kernel matrix, with all
its stabilizer denominators still displayed. The test condition
\texttt{KerEnt.lat\_S} is pairing integrality against
\texttt{KM.nablaX}; its exterior version is integrality against
$T_\lambda^\#$, not against a second copy of $T_\lambda$.
When the source's local integrality assumptions hold, the transported
matrix has exactly the same integral elementary divisors. No new
integrality assumption is silently supplied here.

The precise global integral lattice is also identifiable in the source.
\texttt{Columns/Record.lean}, \texttt{stepH} (lines 1286--1287),
uses \texttt{intDel}; \texttt{Columns/CoefCol.lean},
\texttt{roundPk} (lines 59--64), uses \texttt{Del=delAtr} and
\texttt{Nab=nabAtr}. Thus the integral global column forms use the
Weyl lattice; the dual Weyl lattice supplies its testing lattice.
In \texttt{Divisibility/CoefMod.lean},
\texttt{transLat,delP,delK} (lines 196--211), the lattice is explicitly
the intersection of the inverse images of standard local lattices
under \texttt{actP} applied to $B^{-1}$ for local frames $B$.
Since \texttt{actP} is determinant twist times the transposed
polynomial row action (lines 69--76), its exterior action is
\[
 R_\lambda(B)=\prod_\tau\det(B_\tau)^{s_\tau}
       \bigotimes_{\tau,c}\Lambda^{h_{\tau,c}}B_\tau.
                                                               \tag{IF-08j}
\]
The round realization uses the source's inverse-transpose convention
\texttt{rhoVr}; its \texttt{delAtr} is formed from the dual frame
lattice (\texttt{CoefMod.lean}, lines 320--330). Applying exactly that
same convention to (IF-08j) transports the global intersection and
all its localizations; no free global basis of a nonprincipal
coefficient ideal is presumed.

In the following formula $B_{j,c}$ means the frame of the coefficient
lattice used by \texttt{transLat}. For an actual round column whose
geometric frame is $C_{j,c}$, this means $B_{j,c}=C_{j,c}^{-t}$,
with the dual-frame convention just specified; one must not insert
$C_{j,c}$ directly into $R_j(B_{j,c})$. For chosen local bases at a
prime $w$, the transported block matrix is therefore the explicit matrix
\[
 R_{j+1}(B_{j+1,c'})^{-1}
 \frac{\rho_j(c)\rho'_j(c')}{|\Gamma_{j,c}|}
 \Bigl(\sum_{\beta}J_{j,c',c}(\beta)A_{\beta}\Bigr)
 D_{n_j} R_j(B_{j,c}).                              \tag{IF-08k}
\]
One sums these blocks over the actual finite class representatives,
then restricts to the invariant and level submodules defining forms.
This is the double-coset sum \texttt{kerOpM}; its scalar analogue is
\texttt{thetaLiftAMF} in \texttt{ThetaLiftIndependence/Statement.lean},
lines 80--84. In a local field the scalar valuation in (IF-08k) is
\[
 v_w(\rho_j(c))+v_w(\rho'_j(c'))-v_w(|\Gamma_{j,c}|),
\]
and a jet additionally contributes the recorded period-generator
valuation $\sum_\tau |\beta_\tau|v_w(c_{D,\tau})$ to its ordinary
Taylor derivative valuation. They have not been dropped or replaced
by root factorials. For a resulting local matrix $M$ of rank $r$,
put $s_0=0$ and
\[
 s_k=\min_{|I|=|J|=k}v_w(\det M_{I,J}),\qquad 1\leq k\leq r.
\]
Its elementary-divisor exponents are $s_k-s_{k-1}$; its image index
in the saturated image, when $M$ is integral, has length $s_r$.
These statements follow directly from the Smith normal form over a
DVR, and the formulas are unchanged by (IF-08c). This gives an exact
finite algorithm for the actual global transport, including every
local denominator. A terminal period functional is transported by
the same dual isomorphism, so its value on the transported chain is
identical, not merely bounded by the original value.

The result does not identify (IF-08i) with a single $L$, $L^{(r)}$,
or a polynomial in one Lefschetz pair. The actual moments,
class numbers, characters and stabilizers remain part of the matrix.
It constructs the requested finite-kernel intertwiner and proves
index transport; equality with a separately defined Lefschetz order
index would still require an additional operator identity. IF-03 and
IF-05 show why discarding the operator and lattice distinctions gives
false numerical equalities.

\subsection{IF-09: the actual two-step column chain and a fixed kernel test}
The polynomial coefficient weights of the definite column chain must
be distinguished from the auxiliary $a^+$ row weights in IF-05.
For the actual \texttt{numEx} datum and its all-false word, define
\[
 u=(113/2,47/2,39/2,31/2,23/2,15/2,7/2,-1/2),\qquad
 \lambda_0=(53,21,18,15,12,9,6,3).
\]
The unshifted weights from \texttt{MemberForms.lam\_eq} are
\[
 \lambda_1^N=(\lambda_0,-88^{[16]}),\qquad
 \lambda_2^N=(\lambda_0,-88^{[16]},-90^{[16]}).
\]
The actual \texttt{colVal} uses \texttt{Wts.lamTrue}, which subtracts
$8j$ from every entry at stage $j$. Thus the determinant exponents of
its three polynomial realizations are
\[
 (s_0,s_1,s_2)=(3,-96,-106),
\]
and their polynomial partitions are
\[
\begin{split}
 \alpha_0&=(50,18,15,12,9,6,3,0),\\
 \alpha_1&=(141,109,106,103,100,97,94,91,0^{[16]}),\\
 \alpha_2&=(143,111,108,105,102,99,96,93,2^{[16]},0^{[16]}).
\end{split}                                                     \tag{IF-09a}
\]
Here $x^{[r]}$ means $r$ repeated entries. These formulas follow by
substitution in \texttt{DataCoeff.hcToWeightQ},
\texttt{Params.wG}, \texttt{MemberForms.lam\_eq}, and
\texttt{Columns/AbelianHostFieldConditions.lean},
\texttt{Wts.lamN, Wts.lamTrue} (lines 462--465).
\texttt{Columns/CoefCol.lean}, \texttt{colVal} (lines 89--91), and
\texttt{Divisibility/ColPins.lean}, \texttt{lamOfCol} (lines 21--26),
pin this distinction to the actual column map.

The kernel diagrams, before appending target zeros, are
\[
 \nu_1=(141,109,106,103,100,97,94,91),\qquad
 \nu_2=(143,111,108,105,102,99,96,93,2^{[16]}).       \tag{IF-09b}
\]
They are exactly the output of \texttt{NumDatum.nu} and
\texttt{DV.nuRow}; they are also
$\alpha_0+91(1^8)$ and $\alpha_1+2(1^{24})$.
The source has
$\mathrm{sSharp}=(88,98)$ and $\mathrm{sDiv}=(96,106)$.
The identities
$s_i+\mathrm{sSharp}_i=n_i$ and
$s_{i+1}=-\mathrm{sDiv}_i$ give
\[
 (n_0,n_1)=(91,2).                                  \tag{IF-09c}
\]
They are precisely \texttt{KerEntTied.WtTie} in
\texttt{KernelLattices/JetFieldLeKR.lean}, lines 108--111.
Consequently the actual finite coefficient chain of IF-08 is
\[
 T_{\alpha_0}(8)\xrightarrow{D_{91}}
 T_{\nu_1}(8)\xrightarrow{B_{J_1}}
 T_{\alpha_1}(24)\xrightarrow{D_2}
 T_{\nu_2}(24)\xrightarrow{B_{J_2}}
 T_{\alpha_2}(40),                                  \tag{IF-09d}
\]
The two determinant degree increments are $8\cdot91=728$ and
$24\cdot2=48$. Thus two lifts do not represent $L^2$.
The chain retains the determinant characters $s_i$ and the scalar block
factors (IF-08i). This specifies the whole finite polynomial kernel,
not just a simple-root summand.

The actual column simple-root gaps include $32,3,91,2$ (zero gaps are
trivial). Their top-to-second-descendant invariants are
\[
\begin{array}{c|c|c|c}
 d&F^2\text{ on Weyl}&F^2\text{ on normalized dual}
       &F^{(2)}\text{ on normalized dual}\\\hline
 2&2&2&1\\
 3&2&6&3\\
 32&2&992&496=2^4\cdot31\\
 91&2&8190&4095=3^2\cdot5\cdot7\cdot13
\end{array}.                                                   \tag{IF-09e}
\]
On the Weyl lattice every displayed $F^{(2)}$ invariant is one, by
IF-04. Thus the dual-lattice obstruction is present in the actual
column coefficients themselves. The auxiliary row gap $90$ in IF-05
is a separate, correctly labeled test.

The source's \texttt{numEx} is a numerical datum, not a full
\texttt{NumFacts} witness. The repaired \texttt{numBig} is given in
\texttt{NumbersLink.lean}, \texttt{numBig} (line 803) and
\texttt{numFactsBig} (line 1260). It retains the row parameters and
changes $u_7$ from $-1/2$ to $-7/2$ at one specified real place, besides
changing the size parameters. At that place $s_0=0$, $n_0=88$;
the last entries $91,93$ in (IF-09b) become $88,90$ respectively.
The other entries and the second determinant insertion $n_1=2$
are unchanged. The additional column gaps are $6,88$, with
$\binom62=15$ and $\binom{88}2=3828=2^2\cdot3\cdot11\cdot29$.
The normalized dual obstruction therefore persists for this valid
numerical-facts instance. It does not depend on treating
\texttt{numEx} as a nonexistent full global witness.

There is also a precise host obstruction to regarding
\texttt{numBig} itself as a full global \texttt{D7} record.
\texttt{NumbersLink.numBig} retains \texttt{numEx.H};
\texttt{NumbersLemmas.hostEx\_q,hostEx\_d} give $q=11,d=7$.
But \texttt{DataHost.lean}, \texttt{Host.qd\_big} (line 52), requires
$4tr+1<q,d$. Here $t=15,r=9$, so $4tr+1=541$, contradicting both
inequalities. Thus its \texttt{NumFacts} theorem certifies numerical
conditions only; it cannot instantiate that full global host with
these primes. The polynomial representations and kernel tests above
are nevertheless fully specified and valid on their stated coefficient
lattices.

\paragraph{An exact nonzero kernel test.}
For the rectangular coordinate inclusion $A:\mathbb Z^a\hookrightarrow
\mathbb Z^b$, define
\[
 \Theta_A(Z)=\exp\bigl(\sum_{u,i}A_{ui}z_{ui}\bigr),\qquad
 [Z^\beta]\Theta_A=A^\beta/\beta!.
\]
This is the precise series \texttt{ThE} in
\texttt{SpanPowerMembership/NZRowV.lean}, lines 77--79.
Its factorial cancellation is the supplied theorem
\texttt{fac\_ThE}, lines 89--97; its kernel evaluation is
\texttt{eval\_applyAt0,eval\_Gxz}, lines 99--146.
For this specified series $J(\beta)=A^\beta$, so
$B_J=A_\lambda(A)$ by (IF-08a). Each map in (IF-09d) is then an
inclusion with saturated image. Indeed it carries its ambient tensor
basis injectively into the larger ambient tensor basis; saturation of
the cyclic lattices was proved in IF-08. The determinant insertion
has the same property, since it tensors with a primitive norm-one
vector. Their composite is primitive for the same ambient basis
reason. Its nonzero Smith factors are all one.

Taking the highest vector at the beginning and the normalized dual
highest functional at the end gives period exactly $1$ after both
arrows. This assertion is fully fixed-data: the kernels, lattices,
vectors, determinant insertions and functional have been specified.
On the highest diagonal monomial the derivative factor is
$\prod_i\nu_{j,i}!$ and the coefficient of $\Theta_A$ is its inverse,
so the unit value is an explicit factorial cancellation, not an
assumption that the two normalizations agree.

\paragraph{A compatible common Lefschetz action on this fixed family.}
There is nevertheless a positive integral intertwining result for the
same specified kernel chain. Fix one simple root inside the first eight
rows, for example rows $2,3$. Every determinant insertion and coordinate
inclusion in (IF-09d) commutes with both of its root operators $E_i,F_i$:
$E_i$ and $F_i$ kill a full determinant vector, and the coordinate
inclusion is the identity on those two rows. These identities hold on
the entire cyclic coefficient lattice, not just the highest root string.

For each coefficient lattice apply the particle-hole embedding of IF-04
into $V_{N_j}$, with $N_j=(r_j-1)q_j$ and $q_j$ its number of columns.
Let $N=\max_jN_j$, and pad the embedding with a singleton $e$ in each
of the $N-N_j$ extra symplectic pairs. Singleton padding is an integral
isometry with primitive image and is annihilated by both the raising
and lowering on each extra pair. It therefore intertwines the same
$L,D$ on $V_N$. Primitivity of $T_\lambda$ in its column ambient lattice
(IF-08), followed by the coordinate-basis particle-hole embedding and
padding, makes each image a primitive $L,D$-stable sublattice of $V_N$.
Thus the fixed \texttt{ThE} arrows are specified integral isometric maps
between those sublattices, commuting with one common Lefschetz action,
and their nonzero Smith factors remain one. For (IF-09d), one may take
$N=39\cdot143=5577$ at each embedding factor. Tensoring the eight
factors gives the diagonal root action on $V_{8N}=V_{44616}$, with
the same integral and primitivity statements.

The exterior degree of a vector of simple-root weight $h$ is
$N_j-h$ before padding and $N-h$ after padding. Consequently these
intertwining maps have degree zero. They are not applications of $L$;
that stronger assertion is excluded below. No integral equivariant
extension of the arrows to all of $V_N$ is asserted. Nor is algebraicity
of these particular particle-hole image sublattices asserted; IF-10
provides a separate geometric coefficient realization with its stated
character normalization.

For an arbitrary moment block $B$ between the same standard coefficient
lattices, the exact additional common-root conditions are the finite
matrix equations
\[
 E_i' B=B E_i,\qquad F_i' B=B F_i.                 \tag{IF-09f}
\]
They suffice for all divided powers, by equality in characteristic zero
and their integral preservation. For the actual global form spaces one
must also check that the proposed root action preserves their invariant
and level submodules and their local lattices. These are explicit
operator/lattice conditions after (IF-08k), and are not consequences of
mere theta-kernel rationality. No such common-root invariance is assumed
for the unspecified global moments.

\paragraph{A fixed-normalization obstruction to two ordinary Lefschetz arrows.}
This particular two-step kernel chain cannot be intertwined with two
successive applications of the same fixed ordinary $L$ on one full
integral symplectic exterior lattice by integral embeddings whose
final target image is saturated. Indeed its highest vector is
sent to a primitive highest vector. A saturated integral embedding
preserves primitivity. But $L^2=2L^{(2)}$ on every such exterior
lattice, so the image of $L^2$ is contained in $2V_n$ and contains no
primitive vector. This is a contradiction. Even without saturation,
an isometric integral embedding for the fixed normalized pairings is
impossible: the resulting highest vector has squared norm one, whereas
a nonzero vector in $2V_n$ has squared norm at least four. The full
exterior monomial basis is orthonormal, so both assertions follow
immediately by inspecting its integer coefficients. This proves a
fixed-data failure of the ordinary-$L^2$ normalization, rather than a
failure caused by leaving the test vector unspecified. The same argument
also excludes two positive divided powers of that same fixed $L$:
for $a,b>0$,
\[
 L^{(a)}L^{(b)}=\binom{a+b}{a}L^{(a+b)},
\]
and the integer $\binom{a+b}{a}>1$ divides every coordinate of the
composite image. Thus a primitive or norm-one composite cannot occur.
This leaves tensor-functor realizations such as IF-09d, or operators
belonging to different specified Lefschetz pairs, outside the
obstruction.

These kernels are a concrete algebraic test family for the exact
source operator formula; the source also builds \texttt{KerEnt}
instances with \texttt{ThE} in \texttt{NZRowW.lean},
\texttt{levelRecordTwistEntries}, and \texttt{NZRowZ.lean},
\texttt{fin\_ent}. We do not assert that this exponential test family
is the globally automorphic theta function in Alp\"oge's construction,
or that it satisfies all the additional global theta hypotheses.
For that actual global function, (IF-08i) is still an explicit
intertwining formula in its actual finite jets; computing its Smith
numbers requires those finite jet values and class data. An
existential theta construction or a lower bound does not determine
those numbers. The finite-kernel realization has therefore been
constructed, while the distinct claim that every actual global index
equals one universal Lefschetz order index has not been proved.

\paragraph{Schemes, $\lambda$-operations, and realizations.}
Every $A_\lambda(Z)$ above is a polynomial matrix over
$\operatorname{Spec}\mathbb Z$, assembled from the exterior-power
operations $\Lambda^{h_c}$ and tensor products. Its class in $K_0$
is the corresponding product of $\lambda^{h_c}$ classes. The
saturated cyclic lattice and its explicit dual retain information
which $K_0$ alone forgets: a finite-index inclusion has the same
rational character and can have different integral cokernel.
Base change to $\mathbb Z_\ell$ preserves the displayed integral
intertwiners and carries a Smith factor $d$ to local length
$v_\ell(d)$. Tensoring with $\mathbb Q_\ell$ forgets that length.
When these spaces are realized in the cohomology of an abelian
variety, their tensor and exterior operations are motivic tensor
operations. The formal coefficient identification alone supplies no
new algebraic correspondence representing the theta moment matrix
or the Hodge sign operator. Such an assertion would require the
separate algebraicity comparison under discussion in Q6g and the
main question.


% IF-10: explicit geometric coefficient realization, with every normalization retained.
\subsection{IF-10: a geometric realization over a CM elliptic curve}
There is a geometric realization of the finite coefficient comparison,
but its integral and character normalizations must be stated. It is not
an identification of an automorphic motive with the cohomology of one
abelian variety, and does not replace the operators of IF-08 by one $L$.

Let $K\subset\mathbb C$ be a number field containing the entries of the
finite matrices under consideration. Enlarge it to a number field
$K'$ containing $i$, $K$, and their complex conjugates, and write
$R=\mathcal O_{K'}$. Denominators of the matrices remain allowed and are
recorded by fractional $R$-lattices, as in (IF-08k). Use the elliptic curve
\[
 E:y^2=x^3-x\quad\text{over }\mathbb Q(i),\qquad
 [i](x,y)=(-x,iy).
\]
We work with algebraic correspondences modulo homological equivalence,
with coefficients in $K'$, and their Betti or $\ell$-adic realizations.
All correspondences constructed below are explicit combinations of
coordinate maps, the origin, $[i]$, and linear combinations of their
graphs. No Hodge-conjecture input is needed for this coefficient model.

\paragraph{Statement IF-10.}
For every finite chain of the exterior coefficient spaces and operators
in IF-08 whose normalized matrix entries lie in $K$, there is a chain
of homological motives generated by powers of $E$ and rank-one CM
character twists with the following properties.

\begin{enumerate}
\item Its realization has exactly the matrices of IF-08, including the
moment coefficients, determinant insertions, class sums, stabilizers,
local frames, invariant subspaces and terminal functionals. The
coefficient lattices, with their specified duals, transport isomorphically.
Before coefficient extension every Smith invariant of the same specified
map is identical. After extension, determinantal ideals are their extended
ideals; local valuations change by the ramification index, as specified
below. No equality of absolute finite indices across a field extension
is asserted.
\item Before the character twists, the polynomial coefficient space of
partition $\lambda$ is a submotive of a power of $E$ whose realization is
$T_\lambda\otimes M^{|\lambda|}$, where $M$ is the $i$-eigenmotive in
$h^1(E)$. A determinant twist $s$ changes the exponent to
$w=|\lambda|+rs$. The underlying integral polynomial lattice is a
saturated intersection in integral cohomology after extending coefficients
to $R$.
\item The Fock isometry uses $2^{-|\lambda|}$ times the ambient
Hermitian form on that polynomial eigenspace. The CM projector has
denominator $2$, and the dual eigenline introduces the factor $-2i$.
At primes above $2$ the transported abstract dual lattice must therefore
be distinguished from the native saturated conjugate eigenlattice.
No assertion of an integral cycle projector, or of an integral cup
polarization with the normalized Fock Gram matrix, is made.
\end{enumerate}

\begin{proof}
Choose an oriented integral symplectic basis $e,f$ of $H^1(E,\mathbb Z)$
such that
\[
 [i]^*e=-f,\qquad [i]^*f=e,\qquad
 \psi(e,f)=1.
\]
This is the basis $dx,dy$ on the square analytic torus. Thus
$u=e+if$ satisfies $[i]^*u=iu$ and
\[
 U=R u=\ker([i]^*-i)\subset H^1(E,R).
                                                               \tag{IF-10a}
\]
This line is saturated: $u,f$ form an $R$-basis, since their coefficient
matrix in $e,f$ has determinant one. Set
$\pi_0=[\{0\}\times E]$, $\pi_2=[E\times\{0\}]$, and
$\pi_1=\Delta_E-\pi_0-\pi_2$, with the convention that the displayed
correspondences act by pullback followed by pushforward to the second
factor. They project to cohomological degrees $0,2,1$, respectively.
The transpose graph of $[i]$ acts by $[i]^*$. Hence
\[
 p_+=\frac{1}{2}\bigl(\pi_1-i\,[i]^*\pi_1\bigr)
                                                               \tag{IF-10b}
\]
is an algebraic idempotent in the homological category and its image is
$M$. Its realization has image $K'u$ and its intersection with integral
cohomology is precisely $U$. The idempotence follows directly from
$([i]^*)^2=-1$ on $H^1$; homological equivalence makes this calculation
an equality of correspondences in the stated category.

For the usual Hermitian extension of the orthonormal form on $e,f$,
\[
 (u,u)_{\rm Herm}=2,\qquad
 \psi(u,\overline u)=-2i,\qquad
 (u,u)_{\rm bilinear}=0.                         \tag{IF-10c}
\]
These are three different pairings. Define the normalized Hermitian
form on $U$ by dividing its ambient Hermitian form by $2$; then $u$ has
norm one. Equivalently its value on $U$ is
$\psi(x,\overline y)/(-2i)$. This last expression is restricted to
$U$, where it is integral and positive under the chosen complex
embedding; it is not asserted to be an integral pairing on all of
$H^1(E,R)$.

For a column of height $h$, use the actual cohomological space
$H^h(E^r,K')$. For each $h$-element subset
$I=\{i_1<\cdots<i_h\}$ of the $r$ factors, project to $M$ in those
factors and to $h^0(E)$ in the others. Its integral generator is the
ordered cup product $u_{i_1}\cdots u_{i_h}$. These generators occupy
distinct K\"unneth summands and each is primitive, so their direct sum
is a saturated lattice canonically identified with
$\Lambda^h R^r\otimes U^{\otimes h}$. Applying this construction in
separate blocks for all column heights $h_c$ realizes
\[
 A_\lambda(R)=\bigotimes_c\Lambda^{h_c}R^r
 \quad\text{as}\quad
 A_\lambda(R)\otimes U^{\otimes d}
 \subset H^d(E^{rq},R),\qquad d=|\lambda|,
                                                               \tag{IF-10d}
\]
where $q$ is the number of columns. We use actual ordered cup products
and the resulting permutation signs. In particular this is the ordinary
exterior power of $H^1$ in cohomology, not an unqualified categorical
exterior power of the odd motive $h^1(E)$.

Every summand indexed by a tuple of subsets in (IF-10d) is isomorphic
to $M^{\otimes d}$. The isomorphism orders the selected factors and
discards the degree-zero factors. Its inverse inserts the origin and
the same degree-zero factors. These are induced by projections,
coordinate permutations, and the degree projectors. Compose the
projection onto one such summand, its ordered identification with
$M^{\otimes d}$, and the inverse identification for another summand.
The resulting algebraic correspondences $E_{JI}$ act by
\[
 E_{JI}(v_I)=v_J,\qquad E_{JI}(v_{I'})=0\quad(I'\ne I),
 \qquad E_{JI}E_{LK}=\delta_{I,L}E_{JK}.           \tag{IF-10e}
\]
The necessary signs are fixed by the chosen increasing order of the
degree-one factors. This provides all matrix units on the coefficient
space; it also provides rectangular matrix units between two ambient
powers with the same $d$.

Choose an integral basis for the primitive sublattice $T_\lambda$
proved in IF-08 and let $C$ be its matrix in the orthonormal basis of
$A_\lambda$. The rational matrix
\[
 P_T=C(C^tC)^{-1}C^t                              \tag{IF-10f}
\]
is the projector to its rational span. Inserting its entries into
(IF-10e) makes it an algebraic projector in the homological category.
Its integral intersection is $T_\lambda\otimes R$: primitivity over
$\mathbb Z$ gives a free quotient and is preserved by scalar extension.
All denominators in (IF-10f) are explicit divisors of
$\det(C^tC)$, before cancellation. Together with the factors $2$ from
(IF-10b), they record the denominator of this chosen correspondence.
The ambient tensor basis in (IF-10d) has squared norm $2^d$; after
multiplying the restricted Hermitian form by $2^{-d}$ its Gram matrix
on $T_\lambda$ is the one in IF-04. Thus the normalized Fock isometry
is exact, including its pairing normalization.

Every coefficient matrix $A_\beta$ of IF-08, and hence every
$B_J=\sum_\beta J(\beta)A_\beta$ with $J(\beta)\in K'$, is now an
algebraic correspondence by (IF-10e), restricted on each side by
(IF-10f). Its matrix on the chosen lattice is unchanged. This argument
does not assert that the coefficients of the representing cycle are
integral; the matrix on the specified lattice is integral exactly
when the corresponding matrix of IF-08 is integral.

To handle determinant characters and arrows of different total
degrees, note that the determinant line of $U^{\oplus r}$ is
$M^{\otimes r}$ with generator $u^{\otimes r}$, using ordered cup
products. Negative tensor powers use the abstract dual of the chosen integral
line from the outset, including in the determinant twist. The polarization
identifies $M^\vee$ with the conjugate eigenmotive $\overline M(1)$;
the scalar in (IF-10c) specifies its integral normalization. These
rank-one motives are invertible in the homological category: evaluation
and its transpose, divided by their nonzero scalar, are inverse on the
one-dimensional realization, hence inverse there. Consequently the
coefficient motive for $(\lambda,s)$ is
\[
 \mathcal V_{\lambda,s}=T_{\lambda,K'}\otimes M^w,
 \qquad w=d+rs.
\]
Normalize it by the specified inverse CM character:
\[
 \mathcal V_{\lambda,s}\otimes M^{-w}
       \simeq T_{\lambda,K'}\otimes\mathbf1.       \tag{IF-10g}
\]
Use the abstract dual of the chosen integral rank-one lattice for this
evaluation, so that its generator evaluates to one. Formula (IF-10c)
says exactly what this means inside a conjugate geometric realization:
$U^\vee$ is represented by
$(-2i)^{-1}\overline U(1)$. In particular it is not the native
saturated lattice $\overline U(1)$ at primes above $2$. If one instead
uses native integral eigenlines in a pair of opposite powers
$M^w,M^{-w}$, evaluation has image the ideal generated by
$(2i)^{|w|}$, up to a unit and the ordering sign. This factor is retained,
or is removed by precisely the displayed abstract-dual normalization;
there is no unnoticed saturation at $2$. This $|w|$ formula concerns
the already collapsed pair $M^w,M^{-w}$. If $s<0$ and one begins
instead with the raw native tensor of $d$ copies of $U$ and
$r|s|$ copies of $\overline U(1)$, converting the latter to abstract
duals contributes $(-2i)^{-r|s|}$. Collapsing the raw native tensor by
pairings first contributes $(-2i)^{\min(d,r|s|)}$, up to the ordering
unit. These factors must be retained in any comparison with that
different native presentation.

After (IF-10g), every rectangular matrix, including the determinant
insertion and maps between different $w$, is an algebraic correspondence
between sums of the unit motive. This realizes exactly the original
coefficient matrices because the evaluation generators were fixed to
one. Finite direct sums, invariant subspaces and level subspaces are
handled by matrix projectors over $K'$ and their specified integral
intersections. An arbitrary projective global lattice is retained as
that lattice; a free global basis is not assumed. Thus the block sums
and local-frame conjugations of (IF-08k), and the terminal dual
functional, are unchanged.

Finally a finite torsion-free extension of rings of integers is
faithfully flat. The minors defining every determinantal ideal are
the same minors after extension to $R$; contracting the extended ideal
recovers the original one. Over a local extension of discrete valuation
rings, a valuation is multiplied by the ramification index, and this
known factor recovers the original length. This proves the index
assertion at every prime, with the abstract-dual and cycle-denominator
conventions just specified.
\end{proof}

\paragraph{The actual algebraicity input.}
The coefficient-field hypothesis is already part of a supplied
algebraic theta-chain datum, and is not proved by calling an arbitrary
complex derivative algebraic. The exact interfaces are
\texttt{ThDistC.lean}, \texttt{Amb.JetsRat} (lines 240--242), and
\texttt{KernelLocalFactors/ThDistC.lean}, \texttt{jetK\_spec}
(lines 615--621), which takes that rationality hypothesis and identifies
the descended moment with the period-normalized derivative.
For the actual tied kernel, \texttt{JetFieldLeKR.lean},
\texttt{KerEntTied.KerJets} (lines 113--116), includes the product
\texttt{rhoDet*rhoDet'} in the algebraic coefficient. One therefore
uses the combined moment $\widetilde J=\rho\rho'J$ when only that
combination is certified algebraic; algebraicity of its separate
factors is unnecessary. Alternatively,
\texttt{DataForms.lean}, \texttt{KRat,KL,KL\_spec} (lines 252--259),
gives the algebraic global map together with its exact complex
realization. These declarations are hypotheses and conditional descent
theorems; they do not by themselves construct a new closed global
instance. IF-10 applies to every instance supplied with this data. In
particular the fixed \texttt{ThE} test of IF-09 is unconditional over
$\mathbb Q$ at the level of its finite moments.

\paragraph{Actual weights and the limit of Tate normalization.}
For one embedding factor $\tau$ of the three actual column coefficients
in IF-09, direct summation gives
\[
 (|\alpha_0|,|\alpha_1|,|\alpha_2|)=(113,841,889),
 \qquad (w_0,w_1,w_2)=(137,-1463,-3351).           \tag{IF-10h}
\]
At the modified real place of \texttt{numBig}, all three $w_j$ decrease
by $3$; their successive differences remain $-1600,-1888$.
The actual coefficient module is a tensor product over $\tau$; apply
(IF-10g) to each factor and then tensor. If the same $M$ is used at
every place, its total character exponent is $\sum_\tau w_{j,\tau}$,
not a single entry of (IF-10h). The CM twists in (IF-10g) are therefore
substantial and are not pure
Tate twists. Indeed, under the selected embedding $M^w$ has Hodge type
$(w,0)$. A nonzero algebraic correspondence between
$M^w(a)$ and $M^{w'}(b)$ would preserve Hodge type, requiring
\[
 (w-a,-a)=(w'-b,-b).
\]
It follows that $a=b$ and $w=w'$. Hence in this natural single-eigenline
model the nonzero arrows of the actual coefficient chain cannot become
algebraic by Tate twists alone. The explicit CM-character normalization
is necessary for this construction. This is an obstruction to this
specified model, not a claim about every conceivable realization.

On an $\ell$-adic realization, if the eigenline has character $\eta$,
the untwisted coefficient motive has character $\eta^w$ on its rank-one
factor (with the sum of exponents when tensoring over $\tau$). Formula
(IF-10g) cancels that character. Thus this construction
does not preserve an unspecified automorphic $L$-function or supply its
motivic realization. It does give algebraic correspondences realizing
the fully specified finite coefficient matrices and their exact local
indices. In the language of cote 16, scans 014--018, it constructs
operators in a geometric correspondence ring after the stated
normalizations; IF-09 prevents replacing all such two-step kernels by
positive powers of one fixed Lefschetz element. The $\lambda$-operations
occur in their actual cohomological exterior powers, while $K_0$ alone
forgets the retained lattice indices and the prime-$2$ dual-lattice
factor.

The equation for $E$ and these coordinate correspondences spread over
$\mathbb Z[i,1/2]$. Statements about the transported lattice at rational
prime $2$ concern integral Betti cohomology or the specified $2$-adic
representation lattice. They do not assert good reduction of this
elliptic curve at $2$, or an integral Chow projector there.


\subsection{IF-11: an inhabited two-step algebraic theta-jet specialization}
The numerical examples in IF-09 are not the only source of coefficient
data. A separate supplied construction gives a genuine theta family
with $m=2$ and all the algebraicity needed by its finite kernels. This
does not assert the additional global size or nonzero-period hypotheses.

All paths in this subsection beginning with \texttt{Development/} are
paths in Alp\"oge's Lean archive. Short paths beginning with \texttt{Proofs/}
or \texttt{Definitions/} below are relative to
\texttt{Development/AbelianVarietyTheory/}. The theorem
\texttt{GlobalTheory.KerT.JetsInhab.exists\_colIn\_jets}, in
\texttt{Development/AbelianVarietyTheory/Proofs/GlobalInterpolation/}
\texttt{ExistsHypJetsAll.lean}, lines 73--115, supplies
\[
 \exists B,X,I,\quad
 \texttt{HasLetters}(B)\ \wedge\ I.\texttt{Ok}\ \wedge\ I.\texttt{JetsIn}.
                                                               \tag{IF-11a}
\]
We use this as a cited theorem. Its proof builds its base from
\texttt{ExistsNatBoundPre.lean}, \texttt{exists\_pre\_two\_words}
(lines 29--132), and then fixes the distinguished word in
\texttt{exists\_pre\_sized} (lines 19--71 of
\texttt{ExistsHypJetsAll.lean}). Its rank datum is
\texttt{Formalism/SharpRestriction/InertJacobiLevels/HostFam.lean},
\texttt{Inhab.Rt 0} (lines 14--29). Thus
\[
 k=8,\quad m=2,\quad t=15,\quad f=8,\qquad
 (N_0,N_1,N_2)=(8,24,40).
\]

\paragraph{Statement IF-11.}
The witness constructed in (IF-11a) has the following coefficient
specialization, constant over its eight embedding factors:
\[
\begin{split}
 u&=(49,45,37,29,21,13,5,-1)/2,\\
 v&=(24,20,16,12,8,4,0),\\
 b&=(1,37),\qquad e=(19,55),\\
 \lambda_0&=(21,20,17,14,11,8,5,3).
\end{split}                                                     \tag{IF-11b}
\]
Its actual column determinant exponents and polynomial parts are
\[
 (s_0,s_1,s_2)=(3,-15,-43),
\]
\[
\begin{split}
 \alpha_0&=(18,17,14,11,8,5,2,0),\\
 \alpha_1&=(28,27,24,21,18,15,12,10,0^{[16]}),\\
 \alpha_2&=(48,47,44,41,38,35,32,30,20^{[16]},0^{[16]}).
\end{split}                                                     \tag{IF-11c}
\]
The two finite kernels have shapes
\[
 \nu_1=(28,27,24,21,18,15,12,10),\qquad
 \nu_2=(48,47,44,41,38,35,32,30,20^{[16]}),          \tag{IF-11d}
\]
with determinant insertions $(10,20)$, of degrees $80,480$ per
embedding. All required normalized moments lie in the finite number
field supplied by (IF-11a). Hence IF-08 and IF-10 apply to this actual
theta-jet family, not only to the exponential test series.

The distinct nonzero simple-root gaps in these three coefficient
lattices are $1,2,3,10,20$. The gap $1$ has no second descendant; for
the others the exact second-descendant invariants are
\[
\begin{array}{c|c|c|c}
 d&F^2\text{ on Weyl}&F^2\text{ on normalized dual}
       &F^{(2)}\text{ on normalized dual}\\\hline
 2&2&2&1\\
 3&2&6&3\\
 10&2&90&45=3^2\cdot5\\
 20&2&380&190=2\cdot5\cdot19
\end{array}.                                                    \tag{IF-11e}
\]
On every nonzero displayed Weyl second-descendant line,
$F^{(2)}$ has invariant one. These are coefficient-root indices,
not the values or Smith invariants of the full global theta composition.

\begin{proof}
The parameter construction
\texttt{Proofs/IwasawaDecomposition/ExistsHarishChandraWeights.lean},
\texttt{exists\_preCh\_of\_prescribed} (lines 12--33), gives
\[
 v_l=(k-2-l)4\Delta,\qquad
 b_j=b_1+2hj,\qquad e_j=b_j+h,
\]
and endpoint values $u_0=u_F+1/2$, $u_{k-1}=u_L+1/2$, with
$u_{l+1}=v_l-1/2-\Delta$ for the distinguished all-true word.
The invocation in \texttt{exists\_pre\_two\_words} has
\[
 \Delta=1,\quad h=t+3=18,\quad u_F=24,\quad u_L=-1,
 \quad b_1=1.
\]
Its distinguished word is $0\in\operatorname{Fin}2$, on which
\texttt{decide}(u=0) is true. These substitutions prove (IF-11b).
The passage to the one-word base keeps this weight family. The field
enlargement retains the preliminary rank and weight data, and the
subsequent change of ambient theta data retains the resulting base $B$. Thus the actual witness in (IF-11a) retains these
parameters; no selection of an unrelated existential witness is used.

Subtracting the rank-eight half-sum of positive roots gives the
displayed $\lambda_0$. From \texttt{Params.wG} one obtains
$w_G=(7,27)$, so the unshifted weights are
\[
 \lambda^N_1=(\lambda_0,-7^{[16]}),\qquad
 \lambda^N_2=(\lambda_0,-7^{[16]},-27^{[16]}).
\]
Subtract $8j$ at stage $j$, as required by
\texttt{Wts.lamTrue}, and then subtract the last entry to obtain
(IF-11c). Here
\[
 \mathrm{sSharp}=(7,35),\qquad
 \mathrm{sDiv}=(15,43),\qquad
 n_j=s_j+\mathrm{sSharp}_j=(10,20).
\]
Adding these $n_j$ to each source polynomial part gives (IF-11d),
which also agrees by direct substitution with \texttt{nuRow}.
The determinant degrees are $8\cdot10=80$ and $24\cdot20=480$.
Successive differences in (IF-11c) give the stated root gaps.
The invariant table is now exactly IF-04, with the specified Weyl
or normalized-dual lattice; no integral identification of those two
lattices is assumed.

It remains to check that the cited algebraicity is exactly the one
used by the kernel. In
\texttt{Definitions/AutomorphicMembership/JetsInIffStepField.lean},
\texttt{ColIn.jetNum} (lines 99--103) is
\[
 J(\beta)=\chi_F^{-1}
       \left(\prod_{\tau,u,i}c_{D,\tau}^{\beta_{\tau,u,i}}\right)
       (\partial^\beta\Theta)(0),                 \tag{IF-11f}
\]
where $\Theta$ is the actual \texttt{thInd} family. This is precisely
\texttt{ThDistC.jetC}, with ordinary derivatives and no division by
$\beta!$. Its \texttt{JetsIn} definition (lines 109--113) asserts
membership in the coefficient field for multidegree
$\sum_r\nu_{j,r}$ at each $\tau$, and
\texttt{JetsIn.exists\_eq,isAlgebraic} (lines 142--150) identifies
these numbers with their field-valued versions.

Every coefficient $A_\beta$ used by the polynomial kernel of shape
$\nu_j$ has exactly this multidegree: each exterior column contributes
its height, so the total is $|\nu_j|$ at that embedding. Consequently
the degree-restricted algebraicity hypothesis proves algebraicity of
every moment that actually occurs in the finite operator. An
all-degree hypothesis is unnecessary. The degrees here are
$155,635$ per embedding, or $1240,5080$ over the eight embeddings.

The field existence is not an assertion about the default-valued
function \texttt{desc}. Its supplied proof is
\texttt{Proofs/IwasawaDecomposition/DegreeEqTotDeg.lean},
\texttt{finAlg\_step} (lines 193--226), followed by
\texttt{exists\_jetsIn} (lines 283--313): a finite algebraic generating
set of the required theta jets is placed in the coefficient field,
using the supplied normalized-vacuum algebraicity theorem
\texttt{GlobalTheory.T.vacuum\_entire\_and\_normalizedJets\_algebraic}
in \texttt{VacuumEntireAndNormalizedJetsAlgebraicStatement.lean},
and theta equivariance and lattice reduction.
\texttt{ExistsColInJetsIn.exists\_bound\_jetsIn} (lines 34--40)
specifies the finite algebraic generating set to adjoin; the
\texttt{BaseExt.exists\_base} call in \texttt{exists\_colIn\_jets}
(lines 85--96) performs that adjunction while retaining the preliminary
rank and weight data.
Finally \texttt{Definitions/GlobalTwisting/JetsAllOfColIn.lean},
\texttt{jetsAll\_of\_colIn} (lines 108--116), identifies these
degree-restricted jets with the actual assembled kernel's
\texttt{JetsAll}. This proves exact compatibility with IF-08 and
discharges IF-10's coefficient-field input for this inhabited family.
The cited theorems are used as supplied results; we have not run a
new Lean kernel check of their proofs.
\end{proof}

\paragraph{A concrete derivative test and a full-module invariant.}
Write $\Delta_h(Z)$ for the leading $h$-row, $h$-column minor.
At one embedding factor the two highest flag polynomials for (IF-11d)
are the explicit polynomials
\[
 P_1=\Delta_1\Bigl(\prod_{h=2}^6\Delta_h^3\Bigr)
            \Delta_7^2\Delta_8^{10},\qquad
 P_2=P_1\Delta_{24}^{20}.                         \tag{IF-11g}
\]
For the full coefficient tensor put
$P_j^{\rm tot}=\prod_\tau P_j(Z_\tau)$. In standard coefficient
frames, before the scalar factors in (IF-08k), the normalized
highest-to-highest block coefficient is the moment functional
$J(P_j^{\rm tot})=\sum_\beta[Z^\beta]P_j^{\rm tot}\,J(\beta)$.
For other frames one uses the complete matrix conjugation (IF-08k);
one does not replace it by a scalar change of this one entry.
To prove the standard-frame formula,
apply the target highest functional to $A_{\nu_j}(Z)t_{\nu_j}$.
Cauchy--Binet selects the leading minors at each embedding, which gives
$P_j^{\rm tot}(Z)$; then contract its coefficients with $J$. Thus
(IF-11g) specifies an actual differential-operator test of degrees
$155,635$ per embedding, or $1240,5080$ in total. For the
coordinate exponential moments it equals one; no such value is
asserted for the actual analytic theta family.

The source coefficient representation has the wrong rank to be a
whole symplectic exterior algebra. The Weyl dimension
formula cited in IF-05 gives
\[
 \operatorname{rank}T_{\alpha_0}(8)
  =2^{28}\cdot3\cdot5^2\cdot7\cdot11^2\cdot19\cdot23.
                                                               \tag{IF-11h}
\]
This follows by multiplying
$\prod_{i<j}(\alpha_{0,i}-\alpha_{0,j}+j-i)/(j-i)$; the supplied
exact prime-exponent computation records each factor. Since a full
$V_n$ has rank $4^n$, even rational isomorphism to a whole $V_n$ is
impossible. The eightfold coefficient tensor retains the odd prime
factors. This does not obstruct the primitive embeddings or tensor
realizations constructed above, and says nothing about dimensions
of global invariant-form subspaces.

The per-embedding CM exponents of IF-10 are now
$(99,-205,-1085)$, with sums $(792,-1640,-8680)$ over all embeddings.
The same Tate-only obstruction and explicit CM-character untwist apply.
For the coordinate exponential family at these partitions, the positive
common-root construction of IF-09 fits inside $V_{39\cdot48}=V_{1872}$
per embedding, and inside $V_{14976}$ after tensoring all eight factors.

\paragraph{Exact scope.}
The input in (IF-11a) certifies a genuine algebraic theta-jet family.
It does not assert a nonzero composed terminal period, positive full
\texttt{sg}, or the full \texttt{SizeData}$_2$ hypotheses.
Indeed the separate source check SF-Q6-02 shows that the named witness
has top endpoint $b_T=-3$, incompatible with that full size condition.
No full step-10 divisibility bound is inferred from (IF-11e).
The movement here is that IF-08 and IF-10 now apply to an explicitly
traced, inhabited source theta family with the exact derivative
normalization, rather than to a hypothetical list of algebraic moments.
The algebraic geometry and $\ell$-adic translation are precisely those
of IF-10: specified coefficient correspondences and lattices, with the
CM character and all local denominators retained.


\subsection{IF-12: genuine CM theta kernels need not commute with the fixed root}
The common-root construction at the end of IF-09 is a positive result
for the specified coordinate-exponential test family. It is not a
root-invariance theorem for every theta jet. We give both a completely
specified CM theta counterexample and an existence statement for the
torsion-translation family of an actual nonzero source CM theta vector.
Neither assertion identifies an individual translate with the prescribed
global automorphic class sum.

\paragraph{The source operator and its normalization.}
For a positive Hermitian form $H$, linear in its first argument, put
\[
 U_x f(z)=f(z+x)\exp\{-\pi H(z,x)-\pi H(x,x)/2\}.       \tag{IF-12a}
\]
This is exactly \texttt{Formalism.L2.Th1.H.opU} in
\texttt{Formalism/LocalToGlobal/TorsionIntegrality/}
\texttt{QuadraticSemicharacters/IsSemiChar.lean}, lines 21--22,
and \texttt{TaylorLogMetric.U} in
\texttt{Formalism/LocalToGlobal/KernelLattices/ColumnData.lean},
lines 101--102. The latter uses
$H(z,x)=\operatorname{area}(\Lambda)^{-1}\sum z_i\overline{x_i}$.
We retain the Gaussian in (IF-12a); replacing this operator by bare
translation would change the jets.

For a polynomial $P=\sum_\beta p_\beta Z^\beta$ write
\[
 J_f(P)=\sum_\beta p_\beta\partial^\beta f(0).
                                                               \tag{IF-12b}
\]
These are ordinary derivatives, with no factorial division. This is
the convention of \texttt{ThDistC.derivAt0} and \texttt{jetC}
(\texttt{ThDistC.lean}, lines 56--61). If $P$ is homogeneous of degree $D$ in one
embedding factor, the source normalized functional is
$\chi_F^{-1}c_D^D J_f(P)$. With several embedding factors it is
$\chi_F^{-1}\prod_\tau c_D(\tau)^{D_\tau}J_f(P)$.
For nonzero $\chi_F,c_D(\tau)$ these common factors cannot turn a
nonzero commutator coefficient into zero. This statement concerns the
fixed multidegree, and does not discard factors between different
degrees.

\paragraph{Statement IF-12.}
Use the genuine first-arrow shapes from IF-11:
\[
 \alpha=(18,17,14,11,8,5,2,0),\qquad
 \nu=(28,27,24,21,18,15,12,10).
\]
Let $D_{10}:T_\alpha(\mathbb Z^8)\longrightarrow
T_\nu(\mathbb Z^8)$ be the determinant insertion, and let
\[
 B_f=\sum_\beta J_f(Z^\beta)A_{\nu,\beta}D_{10}:
 T_\alpha(\mathbb C^8)\longrightarrow T_\nu(\mathbb C^{24})
                                                               \tag{IF-12c}
\]
be the exact transported finite kernel of IF-08. Take $F$ to be the
row root $E_{21}$, on both spaces. There is a genuine CM theta vector
$f$ in the source's \texttt{Th1.H.Th} model, given explicitly below,
such that
\[
                 F B_f-B_f F\ne0.                            \tag{IF-12d}
\]
Its period ratio is $i10^{400}$, and its translation point has order
dividing $12$. It is the unitary translate of the source's normalized
product theta \texttt{F0}, not a formal exponential series. It has
the original semicharacter on the common refined period lattice
$12\Lambda^{192}$.

More generally, let $f_0\ne0$ be any entire theta vector for a full
integral polarized lattice $\Lambda\subset\mathbb C^q$, with the
operator (IF-12a), and let $P\ne0$ be any polynomial. For every
integer $\ell\ge2$ there exist $n\ge0$ and
$x\in\ell^{-n}\Lambda$ such that
\[
                         J_{U_xf_0}(P)\ne0.                  \tag{IF-12e}
\]
The translate retains the original semicharacter on
$\ell^n\Lambda$. Applying this to the polynomial in the proof
below gives (IF-12d) in the torsion-translation family of an actual
nonzero source CM theta vector, even if denominators are restricted
to powers of one chosen allowed odd prime. It does not assert failure
of (IF-12d)'s opposite equality for a particular preassigned translate
or global class sum.

\begin{proof}
First isolate an exact coefficient, including the determinant
insertion. Let $t_\alpha$ be the highest tensor and
$t=D_{10}t_\alpha$. In the target, use the first eight coordinates.
The column-height multiplicities of $\nu$ are
\[
          (a_1,\ldots,a_8)=(1,3,3,3,3,3,2,10).
\]
Thus $r=Ft$ is the single tensor obtained by changing $e_1$ to
$e_2$ in the unique $\bigwedge^1$ factor. All factors of height
at least two contain both $e_1,e_2$ and are killed by $F$.
Consequently $r$ is an ambient unit basis vector and $E_{12}r=t$.
The insertion $D_{10}$ commutes with $F$ because a full determinant
has trivial $\mathfrak{sl}_8$ action. With $\Delta_h(Z)$ the leading
$h$-minor of the $24$ by $8$ matrix $Z$, define
\[
 \begin{split}
 Q(Z)&=\prod_{h=2}^{6}\Delta_h(Z)^3\,
                         \Delta_7(Z)^2\Delta_8(Z)^{10},\\
 P(Z)&=(Z_{11}-Z_{22})Q(Z).
 \end{split}                                                  \tag{IF-12f}
\]
The tensor exterior form is the normalized Fock form by IF-04.
Taking the coefficient of $r$ in the commutator on $t_\alpha$
therefore gives the exact identity
\[
       \langle r,(F B_f-B_fF)t_\alpha\rangle=J_f(P).            \tag{IF-12g}
\]
Indeed the first summand has polynomial coefficient $Z_{11}Q$ and
the second has coefficient $Z_{22}Q$. No scalar $c_\nu$ remains:
it has already been included in the normalized Fock/exterior
identification. The functional $\langle r,-\rangle$ is an integral
dual-Weyl test, since $r$ is an integral unit tensor. The polynomial
$P$ has degree $D=155$ and is nonzero.

We construct an explicit genuine theta vector. Put $T=10^{400}$,
$\tau=iT$, $\Lambda_T=\mathbb Z+iT\mathbb Z$, and
\[
 \vartheta_T(z)=\sum_{n\in\mathbb Z}
                   e^{-\pi T n^2}e^{2\pi i n z},\qquad
 g_T(z)=\frac{e^{\pi z^2/(2T)}\vartheta_T(z)}{\vartheta_T(0)}.
                                                               \tag{IF-12h}
\]
Here $\vartheta_T(0)>0$. Let $L_T$ be the period pair with lattice
$\Lambda_T$ and let $\omega_0=-(1+\tau)$. Then
\[
                g_T=\texttt{Th1.F.th0}(L_T,\omega_0/2).
                                                               \tag{IF-12i}
\]
For completeness, the exact source definition is
$\texttt{th0}(L,h)=U_h\operatorname{sgStar}(L)/
(U_h\operatorname{sgStar}(L))(0)$ in
\texttt{EisensteinTranslationProducts/SemiChar.lean}, line 34.
The supplied identity
\texttt{GlobalTheory.SgStarOddTheta.exp\_mul\_oddTheta\_eq\_deriv\_mul\_sgStar}
in \texttt{Development/AbelianVarietyTheory/Proofs/GlobalInterpolation/}
\texttt{GaussianOddJacobiThetaEqDerivMulSgStar.lean}, lines 100--124,
identifies the Gaussian-gauged odd Jacobi theta with
$\operatorname{sgStar}$ times its derivative at zero. Its odd theta
is defined in \texttt{DifferentiableOddTheta.lean}, lines 19--20, by
\[
 e^{\pi i\tau/4+\pi i(z+1/2)}
       \vartheta_T\bigl(z+(\tau+1)/2\bigr).
\]
Substitute $h=-(\tau+1)/2$ into (IF-12a). The linear exponential
in $z$ cancels, leaving a nonzero constant times
$e^{\pi z^2/(2T)}\vartheta_T(z)$. Division by its value at zero
gives (IF-12i), including its scalar. The two hypotheses of that
source identity are supplied by \texttt{Th1.F.thetaLine} in
\texttt{EisensteinTranslationProducts/EllipticOnePole.lean}, line 27,
and \texttt{Th1.F.sgStarBasic} in
\texttt{EisensteinTranslationProducts/QuasiPeriodConstUnique.lean},
lines 21--22. We use these supplied results as cited theorems.
Equivalently, their
quasiperiod law and normalized value at zero give the same identity.

Set $q=24\cdot8=192$ and
$G_T(Z)=\prod_{i,j}g_T(Z_{ij})$. Thus $G_T$ is exactly the source
$\texttt{Th1.H.F0}(L_T,\omega_0,\operatorname{Fin}24\times
\operatorname{Fin}8)$, by
\texttt{QuadraticSemicharacters/ProductThetaMem.lean}, lines 10--11.
Let $a$ be the $24$ by $8$ real matrix with
\[
 a_{11}=1/4,\quad a_{22}=1/3,\quad
 a_{jj}=1/4\ (3\le j\le8),\quad a_{ij}=0\ (i\ne j),
                                                               \tag{IF-12j}
\]
and put $f_T=U_{iTa}G_T$, with
$H(Z,W)=T^{-1}\sum Z_{ij}\overline{W_{ij}}$.
Direct substitution into (IF-12a), without a gauge change, gives
\[
 f_T(Z)=\vartheta_T(0)^{-q}e^{\pi\sum Z_{ij}^2/(2T)}
      \sum_{n\in\mathbb Z^q}
       e^{-\pi T\|n+a\|^2}e^{2\pi i\langle n+a,Z\rangle}.
                                                               \tag{IF-12k}
\]
All these series and their derivatives converge normally on compact
sets: a fixed derivative contributes only a polynomial in $n$,
whereas the Gaussian exponent is negative quadratic in $n$.

This is an allowed analytic theta function in the source's exact
theta model. Its original product theta has semicharacter
$\prod(-1)^{m_{ij}n_{ij}}$ at lattice points $m+iTn$; substitution
in (IF-12h) verifies the law. Since $12a$ is integral, the Heisenberg
commutator for $iTa$ and any vector in $12\Lambda_T^q$ is one.
Consequently $f_T$ has the original semicharacter, here trivial,
on $12\Lambda_T^q$, and belongs to the source's
$\texttt{Th}(H,1,12\Lambda_T^q)$. The latter definition requires
analyticity and the stated translation law, not self-duality of the
refined lattice. Also $iT$ is imaginary quadratic, so the underlying
elliptic curve and its products are CM; $iTa$ is a torsion point of
order dividing $12$. We do not identify this refined lattice with
the restricted self-dual lattice of a particular \texttt{D7.Amb}.

There is a useful exact finite-degree heat formula. Write
$\Delta=\sum_{i,j}\partial_{X_{ij}}^2$. For a homogeneous polynomial
of degree $D$, ordinary differentiation of the Gaussian gives
\[
 \begin{split}
 S_T&:=\vartheta_T(0)^q e^{\pi T\|a\|^2}(2\pi i)^{-D}
                                                    J_{f_T}(P)\\
 &=\sum_{n\in\mathbb Z^q}
    e^{-\pi T(\|n+a\|^2-\|a\|^2)}
       \left(e^{-\Delta/(8\pi T)}P\right)(n+a).
 \end{split}                                                  \tag{IF-12l}
\]
The heat exponential is a finite sum on this polynomial. To verify
its sign and factor, in one variable divide each derivative by
$2\pi i$. Pairing two derivatives on
$e^{\pi z^2/(2T)}$ contributes
$(\pi/T)/(2\pi i)^2=-1/(4\pi T)$; the exponential differential
operator therefore has coefficient $-1/(8\pi T)$. The multinomial
expansion proves (IF-12l) in all variables. This also verifies that
there is no missing multi-index factorial.

We now give an explicit error bound; no numerical evaluation of an
exponentially small theta value is used. Let $M=\|P\|_1$ be its
coefficient $\ell^1$ norm. Expansion of determinants gives
\[
 M\le 2\prod_{h=2}^{6}(h!)^3(7!)^2(8!)^{10}
                  \le2(8!)^{27}<2\cdot10^{135}.                \tag{IF-12m}
\]
At the diagonal matrix (IF-12j),
\[
           P(a)=-4^{-128}3^{-28},\qquad |P(a)|>10^{-156}.
                                                               \tag{IF-12n}
\]
For $R=\max(1,\|x\|_\infty)$ and
$u=qD^2/(8\pi T)$, termwise differentiation gives
\[
 |(e^{-\Delta/(8\pi T)}P)(x)|\le M R^D e^u,
 \quad |(e^{-\Delta/(8\pi T)}P)(a)-P(a)|\le M(e^u-1).
\]
For example, at each differentiation step a coefficient is
multiplied by at most $D^2$ for each of the $q$ summands of $\Delta$;
the degree decreases, so the displayed bound applies to every term.
Here $qD^2=4{,}612{,}800<10^7$, so $u<10^{-393}$ and
\[
                  M(e^u-1)<4\cdot10^{-258}.                   \tag{IF-12o}
\]
For $n\ne0$, using $|a_{ij}|\le1/3$ and integer coordinates gives
\[
 \|n+a\|^2-\|a\|^2\ge\|n\|^2/3,
 \qquad
 (1+\|n\|_\infty)^D\le e^{D\|n\|^2}.
\]
Since $T\ge2D$ and $\pi\ge3$, the absolute tail in (IF-12l) is at
most
\[
 2M\sum_{n\ne0}e^{-T\|n\|^2/2}
 \le16qM e^{-T/2}
 \le128qM/T^2
 <5\cdot10^{-661}.                                            \tag{IF-12p}
\]
To check the middle inequality, set $r=e^{-T/2}$. The
one-coordinate theta sum is at most $1+4r$; thus the sum over
nonzero vectors is at most $(1+4r)^q-1\le8qr$ when $4qr\le1$.
The last bound uses $e^{T/2}\ge T^2/8$. All of these conditions
hold for the specified integer $T$. The bounds (IF-12o)--(IF-12p)
are smaller in total than $|P(a)|/2$, proving $S_T\ne0$ and hence
$J_{f_T}(P)\ne0$. Equations (IF-12g) and (IF-12l) prove the explicit
counterexample. A homogeneous change to CM period coordinates, or
the common source factors in (IF-12b), multiplies this nonzero
coefficient by a nonzero scalar and does not affect the conclusion.

Finally prove (IF-12e), which does not vary the initial CM datum.
For simplicity choose coordinates on $\mathbb C^q$ and write
$H(z,x)=z^{\mathsf t}C\overline x$, with $C$ invertible. If
$d=\deg P$, multiplying $J_{U_x f_0}(P)$ by
$e^{\pi H(x,x)/2}$ gives a polynomial in $\overline x$, of degree
at most $d$, whose coefficients are holomorphic functions of $x$:
\[
 R(x)=\sum_\beta p_\beta\sum_{\gamma\le\beta}
    {\beta\choose\gamma}(-\pi C\overline x)^{\beta-\gamma}
                                      (\partial^\gamma f_0)(x).
                                                               \tag{IF-12q}
\]
Its part of top antiholomorphic degree is
$f_0(x)P_d(-\pi C\overline x)$, where $P_d$ is the leading
homogeneous part of $P$. Since $C$ is invertible and $P_d\ne0$,
choose a multi-index $\eta$ of degree $d$ with nonzero coefficient
$b_\eta$ in $P_d(-\pi C\overline x)$. Wirtinger differentiation
then gives
\[
             \partial_{\overline x}^{\eta}R(x)
                        =\eta!\,b_\eta f_0(x).                \tag{IF-12r}
\]
Every lower antiholomorphic term vanishes, and holomorphic
coefficients have zero antiholomorphic derivatives. Thus $R$ is not
identically zero. Its nonzero set is nonempty and open. A real basis
of the full lattice identifies
$\bigcup_n\ell^{-n}\Lambda$ with a dense subset of
$\mathbb R^{2q}$, so this set contains a point $x$ proving
(IF-12e). If $\lambda\in\ell^n\Lambda$, integrality of
$\operatorname{Im}H$ on $\Lambda$ makes the commutator of
$U_x$ and $U_\lambda$ equal to one. Hence the original theta law
holds on that refined lattice with its original semicharacter.
\end{proof}

\paragraph{Application to the actual eight-factor source family.}
Choose the nonzero source vector
$f_0=B.\texttt{thZero}(8,24)$.
The supplied theorem \texttt{WeilVac.thZero\_ne\_zero}, in
\texttt{Development/AbelianVarietyTheory/Proofs/IwasawaDecomposition/}
\texttt{ExistsWeilOThZeroTriv.lean}, lines 38--42, proves its
nonvanishing under \texttt{B.StdOK}; \texttt{thZero\_mem}
(lines 80--103) places it in the actual CM theta space. Its period
lattice is full: it is built from a nonzero fractional ideal in the
CM field under its CM embeddings, and coordinate tensor powers of
that lattice. The CM embeddings identify its real span with all the
ambient complex coordinates. The source polarization is integral on
this lattice and has positive scalar matrix $\operatorname{area}^{-1}I$.
For a chosen factor $\tau_0$, use
\[
 P_{\mathrm{all}}(Z)=P(Z_{\tau_0})
          \prod_{\tau\ne\tau_0}\prod_{h=1}^{8}
                         \Delta_h(Z_\tau)^{a_h}.               \tag{IF-12s}
\]
It is nonzero and has degree $155$ in each of the eight factors.
The coefficient of the corresponding target tensor in the root
commutator is exactly $J_{U_xf_0}(P_{\mathrm{all}})$.
The last part of the proof therefore supplies a torsion translate
with a nonzero commutator for the actual IF-11 coefficient shapes.
It does not require the CM theta vector to split over embeddings.

For the algebraicity assertion, assume in addition the source
arithmetic hypothesis \texttt{Crv D cc} $\delta$: this consists of
\texttt{IsCurveIdeal} and
\texttt{IsDiscriminantTwelfthRootOfIdealLattice}. The preceding
analytic nonvanishing assertion requires no such arithmetic
hypothesis. Under \texttt{Crv}, the normalized coefficients of
these finitely many jets are algebraic: use
\texttt{thZero\_coefIn} (lines 236--257) and
\texttt{ta\_thMoved} (lines 311--316), in
\texttt{VacuumEntireAndNormalizedJetsAlgebraic.lean}; the latter
combines the former with \texttt{coefIn\_translate} (lines 53--85
of \texttt{CoefInTranslate.lean}). The general
\texttt{TA.transl} (lines 281--290) also transports an already
established \texttt{TA} property under rational translation.
These are properties of the actual normalized CM vector, not a
consequence of the default map \texttt{desc}. A finite number field
contains the finitely many resulting coefficients. Clearing their
denominators gives an integral map with the same nonzero commutator;
no assertion about its Smith factors is needed for this obstruction.

The relation to the actual characteristic kernels is explicit but
one-directional. The source's
\texttt{GlobalTheory.ColNorms.IndPres.xiInd\_expand}, in
\texttt{Proofs/DivisorProofs/CosetIndicatorThetaTorsionTranslates.lean},
lines 102--114, writes each normalized \texttt{xiInd} kernel as a
finite sum of these exact unitary translates. Its
\texttt{pres\_thInd} (lines 182--248) retains the coefficients,
volume factors, character phases, and rational translation points.
Thus their commutators are the same explicit linear combinations of
(IF-12g). This does not say that a nonzero term survives a particular
combination, nor that every torsion translate is an individual
\texttt{thInd} satisfying all fixed global data. In particular,
(IF-12d) refutes automatic common-root invariance of the genuine
theta-translation family; deciding invariance or period valuations
for a prescribed global automorphic block still requires its actual
combination and global invariant lattices from IF-08.

